\documentclass[11pt]{article}
\usepackage[T1]{fontenc}
\usepackage{lmodern,amsmath,amssymb,amsthm,mathtools}
\usepackage[margin=1in]{geometry}
\usepackage{booktabs,microtype,enumitem,xcolor}
\usepackage[numbers,sort&compress]{natbib}
\usepackage{hyperref}
\hypersetup{colorlinks=true,linkcolor=blue!45!black,citecolor=blue!45!black,
urlcolor=blue!45!black,pdftitle={Sharp Berry-Esseen bounds for Bernoulli sums and reductions for general arrays},pdfauthor={Logan Bell}}
\setlength{\emergencystretch}{2em}
\numberwithin{equation}{section}
\newtheorem{theorem}{Theorem}[section]
\newtheorem{proposition}[theorem]{Proposition}
\newtheorem{lemma}[theorem]{Lemma}
\newtheorem{corollary}[theorem]{Corollary}
\theoremstyle{definition}
\newtheorem{problem}[theorem]{Problem}
\theoremstyle{remark}
\newtheorem{remark}[theorem]{Remark}
\DeclareMathOperator{\Var}{Var}
\DeclareMathOperator{\Bern}{Bern}
\DeclareMathOperator{\Bin}{Bin}
\DeclareMathOperator{\supp}{supp}
\DeclareMathOperator{\sgn}{sgn}
\newcommand{\E}{\mathbb E}
\newcommand{\Pp}{\mathbb P}
\newcommand{\R}{\mathbb R}
\newcommand{\N}{\mathbb N}
\newcommand{\Z}{\mathbb Z}
\newcommand{\ind}{\mathbf 1}
\newcommand{\eps}{\varepsilon}
\newcommand{\be}{C_{\mathrm E}}
\newcommand{\ke}{K_{\mathrm E}}
\newcommand{\ph}{\phi}
\title{Sharp Berry--Esseen bounds for Bernoulli sums\\and reductions for general arrays}
\author{Logan Bell}
\date{10 September 2026}
\begin{document}
\maketitle
\begin{abstract}
Independent centered two-point summands with a common support diameter
satisfy the Berry--Esseen inequality with
$C_{\mathrm E}=(3+\sqrt{10})/(6\sqrt{2\pi})$. We prove that this constant is
optimal in the common-diameter class. We also prove the strict inequality
for two centered two-point summands with arbitrary diameters.
The common-diameter proof uses Schulz's binomial theorem, optimization over
success probabilities, and validated finite calculations. Separate results give atom
and short-interval bounds, with an equality condition on the diameters.
For general arrays, we prove a variance-preserving reduction to summands
with at most three support points and an attained quadratic dual formula.
We then reduce the unrestricted sharp bound to two open problems: the bound
for Bernoulli arrays with arbitrary diameters and an inequality for
three-point mixtures.
\end{abstract}

\section{Results}\label{sec:main}
Let $X_1,\ldots,X_n$ be independent centered real random variables with
finite absolute third moments. Throughout, write
\begin{equation}\label{eq:normalization}
 \begin{gathered}
 B^2=\sum_i\sigma_i^2>0,\qquad \sigma_i^2=\E X_i^2,\qquad
 \ell=B^{-3}\sum_i\E|X_i|^3,\\
 \Delta=\sup_{x\in\R}\left|\Pp\{B^{-1}\textstyle\sum_iX_i\le x\}-\Phi(x)\right|.
 \end{gathered}
\end{equation}
Here $\Phi$ and $\phi$ are the standard normal distribution function and
density. The independent Berry--Esseen problem asks for the least universal
constant in $\Delta\le C\ell$. Esseen's lower bound is
\begin{equation}\label{eq:esseen-constant}
 \be=\frac{3+\sqrt{10}}{6\sqrt{2\pi}},\qquad
 \ke=\sqrt{2\pi}\be=\frac{3+\sqrt{10}}6.
\end{equation}
The equality of the unrestricted constant with $\be$ remains open.

Define
\[
 t_*=\sqrt{10}-3,\qquad p_{\mathrm E}=\frac{1-t_*}{2}
 =\frac{4-\sqrt{10}}2,\qquad q_{\mathrm E}=1-p_{\mathrm E}.
\]
The variance-one Esseen law, in the orientation used here, is
\begin{equation}\label{eq:esseen-law}
 \mu_{\mathrm E}
 =p_{\mathrm E}\delta_{\sqrt{q_{\mathrm E}/p_{\mathrm E}}}
  +q_{\mathrm E}\delta_{-\sqrt{p_{\mathrm E}/q_{\mathrm E}}}.
\end{equation}
Its third moment is positive. Esseen's i.i.d. sequence with this law has
$\Delta/\ell\to\be$ \cite{Esseen1956}.

\begin{samepage}
\begin{theorem}[A common support diameter]\label{thm:equal-gap-two-point}
Suppose each $X_i$ has at most two support points and there is an $h>0$
such that every nondegenerate $X_i$ has support diameter exactly $h$.
Then $\Delta\le\be\ell$. The constant $\be$ is optimal over this class.
\end{theorem}
\end{samepage}

\begin{samepage}
\begin{theorem}[Two arbitrary support diameters]\label{thm:two-summands}
For two independent centered random variables, each supported on at most two
points and with positive total variance, $\Delta<\be\ell$.
\end{theorem}
\end{samepage}
Theorem~\ref{thm:two-summands} does not identify the optimal constant for
two summands. For three or more summands, the class with unequal diameters
remains unresolved, even when all diameters are integer multiples of a
common positive number.

The binomial input is Schulz's theorem \cite[Theorem~1]{Schulz2016}:
for $N\sim\Bin(n,p)$, $n\ge1$, $0<p<1$, and $q=1-p$,
\begin{equation}\label{eq:schulz}
 \sup_x\left|\Pp\{(N-np)/\sqrt{npq}\le x\}-\Phi(x)\right|
 <\be\,\frac{p^2+q^2}{\sqrt{npq}}.
\end{equation}
Our common-diameter theorem extends this assertion to unequal success
probabilities. Its proof reduces probability vectors first to two distinct
values, then to a binomial sum and one exceptional Bernoulli variable.
Section~\ref{sec:verification} describes the computational supplement.
Sections~\ref{sec:moments} and~\ref{sec:finite-reduction} give the reductions
for general arrays. The atom and interval estimates in
Section~\ref{sec:atoms} are separate results.

He and Cheng \cite{HeCheng2026} prove the sharp bound for sufficiently
large i.i.d. samples. The companion manuscript \cite{BellSmallLyapunov2026}
adapts their approach to general independent arrays with sufficiently
small Lyapunov ratio. Its proof uses Theorem~\ref{thm:equal-gap-two-point}
to bound the common-diameter Bernoulli array obtained in its local
comparison. The finite results proved here do not require that asymptotic
argument.

For a centered sum $S$ of variance $B^2$, write
\[
 F_S^+(y)=\Pp(S\le y),\qquad F_S^-(y)=\Pp(S<y),\qquad
 H_S^\varepsilon(y)=\be\ell+
 \varepsilon\Phi(y/B)-\varepsilon F_S^\varepsilon(y),
 \quad \varepsilon\in\{+1,-1\}.
\]
The notation $F_S^\varepsilon$ selects the closed CDF for $\varepsilon=+1$
and its left limit for $\varepsilon=-1$. Continuity of $\Phi$ gives
\begin{equation}\label{eq:one-sided}
 \Delta=\sup_y\max\{F_S^+(y)-\Phi(y/B),\ \Phi(y/B)-F_S^-(y)\}.
\end{equation}
For finite support it suffices to inspect these two values at each atom:
between atoms, $F_S^+-\Phi$ decreases and $\Phi-F_S^+$ increases.
Reflection gives $H_S^-(y)=H_{-S}^+(-y)$.

\section{Proof for a common support diameter}\label{sec:equal-gap-two-point}
Let $B_1,\ldots,B_n$ be independent Bernoulli variables with parameters
$p_1,\ldots,p_n$. Put $q_i=1-p_i$ and
\begin{equation}
 N_p:=\sum_{i=1}^nB_i,\qquad
 \mu:=\sum_{i=1}^np_i,\qquad
 V(p):=\sum_{i=1}^np_iq_i,
 \qquad
 \mathcal B(p):=\sum_{i=1}^np_iq_i(p_i^2+q_i^2).
 \label{eq:equal-gap-bernoulli-data}
\end{equation}
Write $F_j(p)=\Pp(N_p\leq j)$, with $F_j=0$ below the support and $F_j=1$
above it. For an integer $k$ and $V(p)>0$, define the closed and strict
deficits
\begin{align}
 \mathcal H_+(p;k)
 &:=\be\frac{\mathcal B(p)}{V(p)^{3/2}}
   +\Phi\!\left(\frac{k-\mu}{\sqrt{V(p)}}\right)-F_k(p),
 \label{eq:equal-gap-closed-deficit}\\
 \mathcal H_-(p;k)
 &:=\be\frac{\mathcal B(p)}{V(p)^{3/2}}
   -\Phi\!\left(\frac{k-\mu}{\sqrt{V(p)}}\right)+F_{k-1}(p).
 \label{eq:equal-gap-strict-deficit}
\end{align}
By~\eqref{eq:one-sided}, the common-diameter bound is equivalent to the
nonnegativity of both deficits for every integer $k$.
Complementation gives the exact reflection
\begin{equation}
 \mathcal H_-(p_1,\ldots,p_n;k)
 =\mathcal H_+(q_1,\ldots,q_n;n-k).
 \label{eq:equal-gap-owner-reflection}
\end{equation}
Thus it is enough, when convenient, to treat the closed CDF value.

We use the following bound for third differences of a Poisson-binomial
mass function.

\begin{samepage}
\begin{lemma}[A bound for third differences]\label{lem:poisson-binomial-third-difference}
Let $c_j=\Pp(C=j)$ be the mass function of any finite Poisson-binomial law,
extended by zero off its support, and let
\[
 W=\Var C,\qquad
 \Delta^3c_j:=c_j-3c_{j-1}+3c_{j-2}-c_{j-3}.
\]
Then
\begin{equation}
 \sup_j|\Delta^3c_j|
 < \frac{8\be}{(W+1)^{3/2}}.
 \label{eq:poisson-binomial-third-difference}
\end{equation}
\end{lemma}
\end{samepage}

\begin{proof}
Set $A=8\be$. The rational estimates $\sqrt{10}>79/25$ and
$\pi<22/7$ give $A>16/5$.

Suppose first that $0\leq W\leq1/4$. Write $C=\sum_iB_i$, where the $B_i$ are
independent with $B_i\sim\Bern(p_i)$, and put $q_i=1-p_i$. Choose the more likely value of each
$B_i$ as a base configuration, put $\varepsilon_i=\min(p_i,q_i)$, and let $J$
count deviations from that configuration. If
$E=\sum_i\varepsilon_i$ and $P_r=\Pp(J=r)$, then
\[
 E\leq2W\leq\frac12,
 \qquad P_0\geq1-E\geq\frac12.
\]
With $R=\sum_i\varepsilon_i/(1-\varepsilon_i)$, one has
$W\leq E\leq R\leq2E\leq1$, $P_1=P_0R$, and
\[
 P_0\geq e^{-R},\qquad
 P_1\geq Re^{-R}\geq We^{-W}.
\]
Let $m$ be the number of ones in the base configuration. At the two central
indices $m+1$ and $m+2$, the no-deviation contribution to $\Delta^3c_j$ has
magnitude three. Every one-deviation contribution has the opposite sign and
magnitude at least one; every remaining contribution has magnitude at most
three. Since $P_0\geq1/2$, the central sign cannot reverse, and hence
\begin{equation}
 |\Delta^3c_j|
 \leq3P_0-P_1+3\Pp(J\geq2)
 =3-4P_1
 \leq3-4We^{-W}
 \label{eq:third-difference-small-central}
\end{equation}
at those two indices. At every other index, the base contribution has
magnitude at most one, so
\begin{equation}
 |\Delta^3c_j|\leq P_0+3(1-P_0)\leq1+4W.
 \label{eq:third-difference-small-outer}
\end{equation}
The outer bound in~\eqref{eq:third-difference-small-outer} is at most two,
whereas $A(1+W)^{-3/2}>2$ on this interval. For the central
bound~\eqref{eq:third-difference-small-central}, use
$e^{-W}\geq1-W$ and
\[
 (1+W)^{-3/2}
 \geq h(W):=1-\frac32W+\frac{15}{8}W^2-\frac{35}{16}W^3.
\]
For
\[
 P(W):=Ah(W)-(3-4W+4W^2),
\]
one has $h'(W)<0$ and, because $A>3$,
\[
 P'(W)<3h'(W)+4-8W
 =-\frac12+\frac{13}{4}W-\frac{315}{16}W^2<0.
\]
At the right endpoint,
\[
 P(1/4)=A\frac{725}{1024}-\frac94>\frac1{64}.
\]
This proves~\eqref{eq:poisson-binomial-third-difference} for $W\leq1/4$.

For $W\geq1/4$, Fourier inversion and
\[
 |1-p+pe^{it}|\leq
 \exp\{-2p(1-p)\sin^2(t/2)\}
\]
give
\begin{equation}
 |\Delta^3c_j|
 \leq I(W):=\frac{16}{\pi}\int_0^{\pi/2}
 \sin^3x\,e^{-2W\sin^2x}\,dx.
 \label{eq:third-difference-fourier-majorant}
\end{equation}
Put
$J_r(W)=\int_0^{\pi/2}\sin^rx\,e^{-2W\sin^2x}\,dx$.
Integration of the derivative of
$\sin^4x\cos x\,e^{-2W\sin^2x}$ yields
\[
 (5+4W)J_5=4J_3+4WJ_7,
 \qquad \frac{J_5}{J_3}\geq\frac4{5+4W}.
\]
Consequently
\[
 \frac{d}{dW}\log\{I(W)(1+W)^{3/2}\}
 =-2\frac{J_5}{J_3}+\frac{3}{2(1+W)}
 \leq-\frac8{5+4W}+\frac3{2(1+W)}<0.
\]
It remains to check $W=1/4$. The bound $e^{-u}\leq1-u+u^2/2$ gives
\[
 I(1/4)
 \leq\frac{16}{\pi}\left(J_3(0)-\frac12J_5(0)+\frac18J_7(0)\right)
 =\frac{256}{35\pi}.
\]
Finally,
\[
 \frac{256}{35\pi}<\frac{A}{(5/4)^{3/2}},
\]
because, after squaring positive quantities, this is
$5760<49\pi(19+6\sqrt{10})$, which follows from
$\pi>157/50$ and $\sqrt{10}>79/25$. The normalized majorant in
\eqref{eq:third-difference-fourier-majorant} is decreasing, completing the
proof.
\end{proof}

\begin{samepage}
\begin{proposition}[Reduction to a binomial sum and one Bernoulli variable]
\label{prop:equal-gap-one-exception-reduction}
Fix $n$, a mean $\mu$, an integer $k$, and either deficit in
\eqref{eq:equal-gap-closed-deficit}--\eqref{eq:equal-gap-strict-deficit}.
If that deficit is negative somewhere on
\[
 \mathcal P_{n,\mu}:=
 \left\{p\in[0,1]^n:\sum_{i=1}^np_i=\mu\right\},
\]
then, after deleting deterministic coordinates and possibly applying the
reflection~\eqref{eq:equal-gap-owner-reflection}, a negative global minimum
has the form
\begin{equation}
 N_p\sim\Bin(M,p)+\Bern(x),
 \qquad M\geq1,\quad 0<p,x<1,\quad p\ne x.
 \label{eq:one-exceptional-family}
\end{equation}
\end{proposition}
\end{samepage}

\begin{proof}
Since $\mathcal B(p)\geq V(p)/2$, extend the selected deficit by $+\infty$
on $V=0$. It then has a global minimum away from the zero-variance boundary.
Choose a face of least dimension on which the selected deficit is negative
somewhere. Among its global minimizers, choose one maximizing
$\sum_i p_i^3$. A coordinate at zero can be deleted. A coordinate at one can be deleted after replacing
$(\mu,k)$ by $(\mu-1,k-1)$. Both operations preserve $k-\mu$ and the selected
deficit. Each face is identified bijectively with the reduced polytope, so
global minimality is preserved; the cubic tie-break is preserved as well,
because the deleted cube is the constant zero or one on that face. Minimal
dimension therefore makes the chosen minimizer interior.

Select three coordinates and write $e_1,e_2,e_3$ for their elementary
symmetric polynomials. Their generating polynomial is
\begin{equation}
 \prod_{i=1}^3\{1+p_i(w-1)\}
 =1+e_1(w-1)+e_2(w-1)^2+e_3(w-1)^3.
 \label{eq:equal-gap-triple-generating}
\end{equation}
Holding $e_1,e_2$ fixed preserves the mean and the triple variance
$V_3=e_1-e_1^2+2e_2$. At every fixed integer threshold, the CDF of the full
sum is affine in $e_3$, and the contribution of the selected triple to
$\mathcal B$ is
\begin{equation}
\begin{aligned}
 \mathcal B_3={}&e_1-3e_1^2+4e_1^3-2e_1^4
 +(6-12e_1+8e_1^2)e_2-4e_2^2\\
 &+4(3-2e_1)e_3.
\end{aligned}
\label{eq:equal-gap-triple-budget}
\end{equation}
Thus the whole deficit is affine on the feasible $e_3$ interval. If the three
roots are distinct and interior, the current $e_3$ is interior. A nonzero
slope permits a decrease of the deficit; a zero slope permits an increase of
$e_3$ and therefore of
\[
 p_1^3+p_2^3+p_3^3=e_1^3-3e_1e_2+3e_3.
\]
Both alternatives contradict the choice of minimizer. Hence it has at most
two distinct interior probabilities.

Schulz's theorem~\eqref{eq:schulz} excludes a negative deficit when all
probabilities are equal.
Apply the reflection
\eqref{eq:equal-gap-owner-reflection} if necessary, so the selected deficit
is closed. Suppose the two values are $0<a<b<1$, with multiplicities $r$ and
$s$. It remains to show that $r,s\geq2$ is impossible.
For a selected triple, let $d_j$ be the mass function of the other
coordinates and put $\Delta^2d_j=d_j-2d_{j-1}+d_{j-2}$. By
\eqref{eq:equal-gap-triple-generating}, the coefficient of $e_3$ in the
selected CDF is $-\Delta^2d_k$. Hence the coefficient of $e_3$ in the
deficit is
\begin{equation}
 L(e_1;d)=4\gamma(3-2e_1)+\Delta^2d_k,
 \qquad \gamma:=\frac{\be}{V^{3/2}}.
 \label{eq:equal-gap-fiber-slope}
\end{equation}

For $f(t)=t^3-e_1t^2+e_2t$, one has
$f''(a)=2(a-b)<0$ at $(a,a,b)$ and $f''(b)=2(b-a)>0$ at $(a,b,b)$.
Thus the values of $e_3$ at these triples are, respectively, the upper and
lower endpoints of the interval for which
$t^3-e_1t^2+e_2t-e_3$ has three roots in $[0,1]$. Global minimality gives
\begin{equation}
 L_{aab}\leq0\leq L_{abb}.
 \label{eq:equal-gap-fiber-endpoint-signs}
\end{equation}
If the other endpoint of this interval first sends a root to zero or one,
the same one-sided endpoint conditions hold; no interior endpoint has been
omitted.

When $r,s\geq2$, delete two copies of each parameter and let $c_j$ be the mass
function of the remaining core. The deletions associated with the two triples
are $c*\Bern(b)$ and $c*\Bern(a)$. Direct substitution in
\eqref{eq:equal-gap-fiber-slope} gives
\begin{equation}
 L_{abb}-L_{aab}
 =(b-a)\{-8\gamma+\Delta^3c_k\}.
 \label{eq:equal-gap-two-block-slope-difference}
\end{equation}
The slope-difference identity~\eqref{eq:equal-gap-two-block-slope-difference}
and the signs in~\eqref{eq:equal-gap-fiber-endpoint-signs} would force
\[
 \Delta^3c_k\geq\frac{8\be}{V^{3/2}}.
\]
If $W$ is the variance of the core, then
$V=W+2a(1-a)+2b(1-b)\leq W+1$. Lemma~\ref{lem:poisson-binomial-third-difference}
gives the opposite inequality. Thus one of $r,s$ equals one, which is exactly
the family~\eqref{eq:one-exceptional-family}.
\end{proof}

For a binomial sum and one Bernoulli variable, three adjacent binomial
masses bound the second derivative of the deficit. The following inequality
will exclude a negative minimum when the binomial sum has at least two
summands. The two-summand case is proved in Section~\ref{sec:two}.

\begin{samepage}
\begin{lemma}[A curvature inequality on a bounded domain]
\label{lem:one-exceptional-compact}
Let $0<p,x<1$, $q=1-p$, and
\[
 pq\leq W\leq3,\qquad 0\leq f\leq1,
 \qquad A_0:=W-(1-f)q^2\geq0.
\]
Set
\[
 B_0:=p^2+(1-f)(W+2pq),\qquad
 V:=x(1-x)+W+pq,
\]
\[
 t_x:=1-2x,\qquad t_p:=1-2p,\qquad
 y:=\frac{f(W+pq)}{1-qf},\qquad z:=\frac{y-x}{\sqrt V}.
\]
Then
\begin{equation}
 \frac{2\be}{\sqrt V}\{4W+2+t_xt_p\}
 +z\ph(z)(2+z^2)
 -5V\frac{fB_0}{(2-f)B_0+A_0}>0.
 \label{eq:one-exceptional-compact-gap}
\end{equation}
\end{lemma}
\end{samepage}

\begin{proof}
For $V\leq1/5$, put $S_2=t_x^2+t_xt_p+t_p^2$. The inequalities
$V\geq\{2-t_x^2-t_p^2\}/4$ and
$S_2\geq(t_x^2+t_p^2)/2$ imply $S_2\geq1/3$.
Since $4V+S_2=4W+2+t_xt_p$, the arithmetic--geometric mean inequality
bounds the first term in~\eqref{eq:one-exceptional-compact-gap} from below by
\[
 \be\left(8\sqrt V+\frac{2S_2}{\sqrt V}\right)
 \geq\frac{8\be}{\sqrt3}.
\]
The normal contribution is at least
$-4/(e\sqrt\pi)$, since $u(2+u^2)\ph(u)$ is maximized at $u=\sqrt2$.
Moreover
\[
 (2-f)B_0+A_0-fB_0=2(1-f)B_0+A_0\geq0,
\]
so the neighbor fraction is at most one and its entire subtracted
contribution is at most $5V\leq1$. The rational bounds
\[
 \be>\frac25,\qquad \sqrt3<\frac{26}{15},\qquad
 e>\frac{19}{7},\qquad \sqrt\pi>\frac74
\]
give the strict margin
\begin{equation}
 \frac{8\be}{\sqrt3}-\frac4{e\sqrt\pi}-1
 >\frac{24}{13}-\frac{16}{19}-1=\frac1{247}.
 \label{eq:one-exceptional-small-variance-margin}
\end{equation}

For the complement $V\geq1/5$, the normal term is bounded below by
\begin{equation}
 z\ph(z)(2+z^2)\geq\sqrt{\frac2\pi}\min\{z,0\}.
 \label{eq:one-exceptional-normal-lower}
\end{equation}
Indeed the claim is immediate for $z\geq0$; for $z=-u<0$, the function
$(2+u^2)\ph(u)$ decreases because its derivative is $-u^3\ph(u)$, and its
value at zero is $2\ph(0)=\sqrt{2/\pi}$. We verify the remaining domain
using 192-bit Arb intervals and the following three parametrizations:
\begin{enumerate}[label=(\roman*),leftmargin=2em]
\item $p=(1+s_p)/2$, $W=pq+s_W(3-pq)$, $f=s_f$;
\item $p=s_p/2$, $W=q^2+s_W(3-q^2)$, $f=s_f$;
\item $p=s_p/2$, $W=pq+s_W(q^2-pq)$,
      $f=1-(1-s_f)W/q^2$,
\end{enumerate}
where $s_p,s_W,s_f,x\in[0,1]$. The first chart is exactly $p\geq1/2$. For
$p\leq1/2$, the other two charts split at $W=q^2$; below that interface,
$A_0\geq0$ is equivalent to $f\geq1-W/q^2$, which is precisely the range of
the third parametrization. The interfaces overlap. Thus the charts cover the
whole stated domain. In the third chart $A_0=s_fW$; the first two
charts give, respectively,
$A_0=qs_p+s_W(3-pq)+fq^2$ and
$A_0=s_W(3-q^2)+fq^2$. On each box the program uses a downward enclosure for
the smooth term, the lower bound~\eqref{eq:one-exceptional-normal-lower} for
the normal term, and an upward enclosure for the subtracted mass term in
\eqref{eq:one-exceptional-compact-gap}. A box is accepted only when the
resulting outward lower endpoint is positive. Otherwise the widest normalized
coordinate is bisected dyadically;
an unresolved box at the prescribed depth or box budget aborts the program.
Boxes touching both $p=0$ and $f=1$ use the weaker limit $y\geq0$ rather than dividing by an
interval containing zero. If interval dependency does not certify the
three-neighbor denominator as positive, the program does not divide; it uses
the universal bound of one for the neighbor fraction and accepts only if the resulting lower gap is
positive. The finite cover and its strict acceptance bounds are recorded in
Section~\ref{sec:verification}.

The verifier is
\path{verify/finite/common_diameter/singleton_certificate.py}.
The supporting records are listed in \path{certificates/finite/}.
\end{proof}

\begin{samepage}
\begin{proposition}[A binomial sum and one Bernoulli variable]
\label{prop:one-exceptional-bernoulli}
For every $M\geq1$, $0<p,x<1$, $p\ne x$, and every integer $k$, both deficits
of $\Bin(M,p)+\Bern(x)$ are nonnegative.
\end{proposition}
\end{samepage}

\begin{proof}
By~\eqref{eq:equal-gap-owner-reflection}, it suffices to treat the closed
CDF value. Suppose it has a negative deficit. Proposition~\ref{prop:equal-gap-one-exception-reduction},
applied with the same mean and threshold, gives a negative global minimum
of the same form with $p\ne x$. Relabel its parameters as $M,p,x$.

First suppose $M\geq2$, put $m=M-1$, and let
\[
 D\sim\Bin(m,p),\qquad d_j=\Pp(D=j),\qquad W=mpq.
\]
Vary the success probabilities of the exceptional variable and one variable
in the repeated group, keeping their sum fixed. Write their probabilities as $\bar p+d$ and $\bar p-d$, and put
$u=d^2$. Their generating polynomial is
\[
 \{1+\bar p(w-1)\}^2-u(w-1)^2,
\]
so the selected closed CDF along the smoothing curve satisfies the exact
affinity
\begin{equation}
 \mathcal F_k(u)=\mathcal F_k(0)+(d_{k-1}-d_k)u.
 \label{eq:one-exceptional-pair-cdf-affinity}
\end{equation}
At the current point set
\begin{align*}
 t_x&=1-2x, & t_p&=1-2p, & \xi&=d_{k-1}-d_k,\\
 S_2&=t_x^2+t_xt_p+t_p^2, &
 z&=\frac{k-x-Mp}{\sqrt V}, & V&=x(1-x)+Mpq.
\end{align*}
The moving pair satisfies
\[
 V'=-2,\qquad \mathcal B'=-2S_2,\qquad \mathcal B''=-8.
\]
If
\[
 \mathcal A(u)=\be\frac{\mathcal B(u)}{V(u)^{3/2}}
 +\Phi\!\left(\frac{k-\mu}{\sqrt{V(u)}}\right),
\]
then $\mathcal H_+=\mathcal A-\mathcal F_k$. Since the current unequal pair
is an interior point of the $u$ interval, stationarity and
\eqref{eq:one-exceptional-pair-cdf-affinity} give $\mathcal A'=\xi$, while
$\mathcal H_+''=\mathcal A''$. Direct differentiation gives
\begin{align*}
 \mathcal A'
 &=\be\frac{\mathcal B'V+3\mathcal B}{V^{5/2}}
   +\frac{z\ph(z)}V,\\
 \mathcal A''
 &=\be\frac{\mathcal B''V^2+6\mathcal B'V+15\mathcal B}{V^{7/2}}
   +\frac{z\ph(z)(3-z^2)}{V^2}.
\end{align*}
Eliminating $\mathcal B$ from these two identities with
$\mathcal A'=\xi$ gives the exact curvature formula
\begin{equation}
 V^2\mathcal H_+''
 =5V\xi
 -\be\left(8\sqrt V+\frac{2S_2}{\sqrt V}\right)
 -z\ph(z)(2+z^2).
 \label{eq:one-exceptional-stationary-curvature}
\end{equation}
A local minimum would require the left side to be nonnegative.

Extend $d_j$ by zero. The adjacent difference is
\begin{equation}
 \xi=
 \begin{cases}
 0, & k<0,\\
 -d_0, & k=0,\\
 d_{k-1}\dfrac{k-(m+1)p}{kq}, & 1\leq k\leq m+1,\\
 0, & k\geq m+2.
 \end{cases}
 \label{eq:one-exceptional-rank-split}
\end{equation}
If $\xi\leq0$, the quantity subtracted from $5V\xi$ in
\eqref{eq:one-exceptional-stationary-curvature} is strictly positive. Indeed
$V\geq x(1-x)+pq$ implies
\[
 \max\{V,S_2\}\geq\frac13,
\]
by the same elementary calculation used in the proof of
Lemma~\ref{lem:one-exceptional-compact},
and hence
\[
 \be\left(8\sqrt V+\frac{2S_2}{\sqrt V}\right)
 +z\ph(z)(2+z^2)
 \geq\frac{8\be}{\sqrt3}-\frac4{e\sqrt\pi}>0.
\]
The right side of~\eqref{eq:one-exceptional-stationary-curvature} is therefore
negative. Thus only the positive range in the third line of
\eqref{eq:one-exceptional-rank-split} remains.

Put $y=k-(m+1)p>0$, $f=y/(kq)$, and $j=k-1$. Then $0<f\leq1$ and the binomial
recurrences give
\[
 d_{j+1}=(1-f)d_j,\qquad
 d_{j-1}=r_0d_j,\qquad
 r_0=\frac{(k-1)q}{(m-k+2)p}.
\]
The mass of these three atoms is at most one, so
\begin{equation}
 \xi=fd_j\leq\frac{f}{2-f+r_0}.
 \label{eq:one-exceptional-three-neighbor}
\end{equation}
Eliminating $m$ and $k$ gives
\begin{equation}
 r_0=\frac{A_0}{B_0},\qquad
 A_0=W-(1-f)q^2,\qquad
 B_0=p^2+(1-f)(W+2pq),
 \label{eq:one-exceptional-neighbor-elimination}
\end{equation}
as well as
\begin{equation}
 y=\frac{f(W+pq)}{1-qf},\qquad
 \be\left(8\sqrt V+\frac{2S_2}{\sqrt V}\right)
 =\frac{2\be}{\sqrt V}\{4W+2+t_xt_p\}.
 \label{eq:one-exceptional-phase-compression}
\end{equation}
Finally, the normal term obeys the lower bound
\eqref{eq:one-exceptional-normal-lower}.
Here $W=mpq\geq pq$, and
\eqref{eq:one-exceptional-neighbor-elimination} gives $A_0\geq0$. For
$W\leq3$, Lemma~\ref{lem:one-exceptional-compact}, together with
\eqref{eq:one-exceptional-three-neighbor}--\eqref{eq:one-exceptional-phase-compression},
makes the right side of~\eqref{eq:one-exceptional-stationary-curvature}
strictly negative, contradicting local minimality.

For $W\geq3$, Fourier inversion gives the analytic bound
\begin{equation}
 |d_{k-1}-d_k|
 \leq\frac4\pi\int_0^{\pi/2}\sin u\,e^{-2W\sin^2u}\,du
 <\frac2{\pi W}.
 \label{eq:one-exceptional-fourier-tail}
\end{equation}
The last inequality follows on setting $s=\cos u$ and using
$1-s^2\geq1-s$. Combining~\eqref{eq:one-exceptional-fourier-tail} with the
curvature formula~\eqref{eq:one-exceptional-stationary-curvature}, and using
$W\leq V\leq W+1/2$, the negative of the right side
of~\eqref{eq:one-exceptional-stationary-curvature} is at least
\[
 8\be\sqrt V-\frac4{e\sqrt\pi}-\frac{10V}{\pi W}.
\]
This expression is concave in $V$, so its minimum is at $V=W$ or
$V=W+1/2$. Both endpoint functions increase with $W$. At $W=3$, the rational
bounds $\be>2/5$, $\sqrt3>5/3$, $\sqrt{7/2}>11/6$, $\pi>3$, and
$4/(e\sqrt\pi)<1$ give the respective lower bounds $1$ and $44/45$.
This excludes $M\geq2$.

When $M=1$, Theorem~\ref{thm:two-summands}, proved independently in
Section~\ref{sec:two}, gives the required bound. Its proof uses the support
endpoint estimates and two direct compact inequalities, and does not use the
common-diameter theorem.
This excludes the assumed negative closed deficit in the remaining case.
Reflection~\eqref{eq:equal-gap-owner-reflection} supplies the left CDF limit.
\end{proof}

\begin{proof}[Proof of Theorem~\ref{thm:equal-gap-two-point}]

Delete deterministic coordinates and scale $h$ to one. Each remaining summand can then be written
$X_i=B_i-p_i$, so the data in~\eqref{eq:equal-gap-bernoulli-data} are exactly
the variance and sum of absolute third moments of the sum. If
the bound in Theorem~\ref{thm:equal-gap-two-point} failed, one of the deficits
\eqref{eq:equal-gap-closed-deficit}--\eqref{eq:equal-gap-strict-deficit}
would be negative at an integer threshold. Proposition~\ref{prop:equal-gap-one-exception-reduction}
would produce a negative minimum in the family~\eqref{eq:one-exceptional-family},
contrary to Proposition~\ref{prop:one-exceptional-bernoulli}.
Esseen's i.i.d.\ Bernoulli sequence has equal support diameters and its ratio
tends to $\be$, so no smaller constant is valid in this class.
\end{proof}


\section{Two summands with arbitrary diameters}\label{sec:two}
For independent Bernoulli variables $B_i$ with $\Pp(B_i=1)=p_i\in(0,1)$,
write $q_i=1-p_i$ and
\[
 X_i=h_i(B_i-p_i),\quad a_i=p_iq_i,\quad r_i=p_i^2+q_i^2,\quad
 V=\sum_i a_i h_i^2=B^2,\quad R=\sum_i a_ir_ih_i^3=B^3\ell.
\]
We first record the estimate used when a compact calculation reaches a
small variance.
\begin{samepage}
\begin{lemma}[Optimization over the diameters]\label{lem:reciprocal-gap-resource}
Let $0<p_i<1$ and $h_i\geq0$, with $\sum_i a_i h_i^2>0$. Then
\begin{equation}\label{eq:reciprocal-gap-resource}
 \frac{\sum_i a_ir_ih_i^3}{(\sum_i a_ih_i^2)^{3/2}}
 \geq\left(\sum_i\frac{a_i}{r_i^2}\right)^{-1/2}.
\end{equation}
Equality holds precisely when $r_i h_i$ is constant in $i$.
\end{lemma}
\end{samepage}
\begin{proof}
H\"older's inequality gives
\[
 \sum_i a_i h_i^2
 =\sum_i(a_ir_ih_i^3)^{2/3}(a_i/r_i^2)^{1/3}
 \leq\left(\sum_i a_ir_ih_i^3\right)^{2/3}
      \left(\sum_i a_i/r_i^2\right)^{1/3}.
\]
Its equality condition is $a_ir_ih_i^3$ proportional to $a_i/r_i^2$.
\end{proof}

The endpoint estimate below holds for any finite number of summands. It is
also used in Section~\ref{sec:finite-reduction}.
\begin{samepage}
\begin{proposition}[Support endpoints]\label{cor:all-support-owners}
For every finite Bernoulli array with $0<p_i<1$ and $h_i>0$, both
$H_S^+(y)$ and $H_S^-(y)$ are strictly positive at each endpoint of the
support of $S=\sum_i h_i(B_i-p_i)$.
\end{proposition}
\end{samepage}
The proof is in Appendix~\ref{app:support}. It reduces the two CDF values
containing an endpoint atom to an explicit rational inequality for the
Bernoulli odds. The other two values follow from an elementary normal-tail
estimate.

The interior atoms of a two-summand law are covered by the following two
compact inequalities.
\begin{samepage}
\begin{lemma}[Two compact inequalities]\label{lem:two-compact}
Let $0<p_1,p_2<1$, $q_i=1-p_i$, $a_i=p_iq_i$, and $r_i=p_i^2+q_i^2$.
For $0\leq t\leq1$, put $V=a_1t^2+a_2$ and $R=a_1r_1t^3+a_2r_2$. Then
\begin{equation}\label{eq:two-upper-middle}
 \be\frac{R}{V^{3/2}}+
 \Phi\left(\frac{-p_1t+q_2}{\sqrt V}\right)-(1-p_1p_2)>0.
\end{equation}
For $0\leq\theta\leq\pi/2$, set $\alpha=\cos\theta$, $\beta=\sin\theta$.
Then
\begin{equation}\label{eq:two-lower-middle}
 \be\left(\alpha^3\frac{r_1}{\sqrt{a_1}}+
                  \beta^3\frac{r_2}{\sqrt{a_2}}\right)
 +\Phi\left(\alpha\sqrt{\frac{q_1}{p_1}}
            -\beta\sqrt{\frac{p_2}{q_2}}\right)-q_2>0.
\end{equation}
\end{lemma}
\end{samepage}
The proof, rational covers, and bounds at singular parameter faces are given
in Section~\ref{sec:verification}. Inequality~\eqref{eq:two-lower-middle}
is proved for every $\alpha,\beta\geq0$ with $\alpha^2+\beta^2=1$.
The parameters obtained from ordered diameters $0<h_1\leq h_2$ form a
subset of this domain.

\begin{proof}[Proof of Theorem~\ref{thm:two-summands}]
Delete any deterministic coordinate. A remaining single Bernoulli variable
is covered by Proposition~\ref{cor:all-support-owners}. Otherwise relabel
and scale so $h_2=1$ and $0<h_1=t\leq1$.
When $t<1$, the four configurations have the order $00<10<01<11$.
If $x_A=\sum_i(A_i-p_i)h_i$, their CDF values are
\[
\begin{array}{c|cc}
 A&F_S^-(x_A)&F_S^+(x_A)\\\hline
 00&0&q_1q_2\\
 10&q_1q_2&q_2\\
 01&q_2&1-p_1p_2\\
 11&1-p_1p_2&1.
\end{array}
\]
Proposition~\ref{cor:all-support-owners} handles both endpoint atoms.
The closed value at $01$ is \eqref{eq:two-upper-middle}.
At $10$, put
\[
 \alpha=\frac{h_1\sqrt{a_1}}{\sqrt V},\qquad
 \beta=\frac{h_2\sqrt{a_2}}{\sqrt V};
 \qquad \alpha^2+\beta^2=1.
\]
Its closed deficit becomes \eqref{eq:two-lower-middle}.
At $t=1$ the two middle atoms merge; their left and closed CDF values are
$q_1q_2$ and $1-p_1p_2$. Thus \eqref{eq:two-upper-middle} is exactly the
merged closed deficit. The split value $q_2$ is not used at this collision.
Complementing both Bernoulli variables gives $-S$ with the same variance,
third absolute moments, and gap ordering. The identity
$F_S^-(y)=1-F_{-S}^+(-y)$ supplies every left limit from the closed values
just proved. Equation~\eqref{eq:one-sided} completes the proof.
A zero diameter or a probability equal to zero or one is handled by deletion
before standardization. If no variance remains, the hypotheses of the theorem
are not satisfied.
\end{proof}

\begin{samepage}
\begin{corollary}[Mixtures of the two Esseen orientations]\label{cor:mixtures}
Fix $c>0$ and numbers $0\leq\eta_i\leq1$. Suppose the independent variables
$X_i$ have laws
\[
 \mathcal L(X_i)=(1-\eta_i)\mathcal L(cZ)+\eta_i\mathcal L(-cZ),
 \qquad Z\sim\mu_{\mathrm E}.
\]
Then $\Delta\leq\be\ell$.
\end{corollary}
\end{samepage}
\begin{proof}
Represent each mixture by an independent choice of orientation, followed by
an independent draw from that oriented law. Conditional on the orientations,
every summand has variance $c^2$, absolute third moment
$c^3 b_{\mathrm E}$, where
$b_{\mathrm E}=(p_{\mathrm E}^2+q_{\mathrm E}^2)/\sqrt{p_{\mathrm E}q_{\mathrm E}}$,
and common support diameter $c/\sqrt{p_{\mathrm E}q_{\mathrm E}}$.
Theorem~\ref{thm:equal-gap-two-point} bounds every conditional CDF by the
same $\be\ell$ relative to the same standard normal CDF. Averaging preserves
that bound. For $0<\eta_i<1$ the summand has four support points.
\end{proof}

\section{Atoms and short intervals}\label{sec:atoms}
We use the product-measure Lubell--Yamamoto--Meshalkin (LYM) inequality in the form stated by Yehuda and
Yehudayoff \cite[Theorem~2]{YehudaYehudayoff2026}; their Theorem~1 gives the
antichain bound used for the atom comparison. The increasing chain with
rank-conditioned Bernoulli marginals was constructed earlier by Broman,
van de Brug, Kager, and Meester \cite[Lemma~2.2]{BromanEtAl2012}; see also
\cite[Lemma~3 and Proposition~5]{YehudaYehudayoff2026}. We give an inductive
construction, derive the equality condition, and extend the argument to
closed intervals.
\begin{samepage}
\begin{proposition}[Weighted antichains and equality in the atom bound]\label{prop:weighted-antichain}
Let $B_1,\ldots,B_n$ be independent Bernoulli variables with parameters
$p_i\in(0,1)$, let $h_i>0$, and put $N=\sum_iB_i$. For every antichain
$\mathcal A\subseteq2^{[n]}$ (no member properly contains another),
\begin{equation}
 \sum_{A\in\mathcal A}
 \frac{\Pp(\{i:B_i=1\}=A)}
      {\Pp(N=|A|)}
 \leq1.                                                \label{eq:weighted-lym}
\end{equation}
Consequently,
\begin{equation}
 \sup_y\Pp\left(\sum_i h_iB_i=y\right)
 \leq\max_k\Pp(N=k).                                   \label{eq:weighted-atom-comparison}
\end{equation}
If the mass function of $N$ has a unique mode $m\in\{1,\ldots,n-1\}$,
then equality in the atom comparison~\eqref{eq:weighted-atom-comparison}
forces $h_1=\cdots=h_n$.
\end{proposition}
\end{samepage}

\begin{proof}
Write $r_i=p_i/(1-p_i)$ and let $\pi_k$ be the conditional law of the random
subset $\{i:B_i=1\}$ given $N=k$. Thus
\[
 \pi_k(A)=\frac{\prod_{i\in A}r_i}{e_k(r_1,\ldots,r_n)}
 \qquad(|A|=k),
\]
where $e_k$ is the elementary symmetric polynomial. We first construct a
random maximal chain $A_0\subset A_1\subset\cdots\subset A_n$ whose rank-$k$
marginal is $\pi_k$.

Proceed by induction on $n$. Write $\pi_j^{(-n)}$ for the analogous rank-$j$
law on $[n-1]$. Let $\alpha_k=\pi_k\{n\in A\}$, with $\alpha_0=0$ and
$\alpha_n=1$. For $1\leq k\leq n-1$,
\[
 \frac{\alpha_k}{1-\alpha_k}
 =r_n\frac{e_{k-1}(r_1,\ldots,r_{n-1})}
          {e_k(r_1,\ldots,r_{n-1})}.
\]
Newton's inequalities make the ratio on the right nondecreasing in $k$, so
$\alpha_k\leq\alpha_{k+1}$. Couple ranks $k$ and $k+1$ in three parts. With
probability $\alpha_k$, both sets contain $n$ and their intersections with
$[n-1]$ use the inductive coupling from $\pi_{k-1}^{(-n)}$ to
$\pi_k^{(-n)}$. With probability $1-\alpha_{k+1}$, neither set contains
$n$ and the intersections use the coupling from $\pi_k^{(-n)}$ to
$\pi_{k+1}^{(-n)}$. With the remaining probability
$\alpha_{k+1}-\alpha_k$, sample $C\sim\pi_k^{(-n)}$ and use
$C\subset C\cup\{n\}$. Zero-mass branches are omitted; thus no nonexistent
rank is invoked when $k=0$ or $k=n-1$. Separating the contain-$n$ and
exclude-$n$ cases verifies the two marginals. Composing these adjacent
kernels gives the required chain.

Every maximal chain meets an antichain at most once. Taking expectations of
the number of intersections gives the weighted LYM inequality~\eqref{eq:weighted-lym}.
For fixed $y$, the fiber
\[
 \mathcal F_y=\left\{A:\sum_{i\in A}h_i=y\right\}
\]
is an antichain because every $h_i$ is positive. If
$s_k=\Pp(N=k)$ and $q_k=\pi_k(\mathcal F_y)$, then
\[
 \Pp\left(\sum_i h_iB_i=y\right)=\sum_ks_kq_k
 \leq\left(\max_ks_k\right)\sum_kq_k
 \leq\max_ks_k,
\]
which proves the atom comparison~\eqref{eq:weighted-atom-comparison}.

Suppose equality holds, let $y$ attain the supremum, and let $m$ be the
unique interior mode. If $M=\max_k s_k$, then the identity
\[
 M-\Pp\left(\sum_i h_iB_i=y\right)
 =M\left(1-\sum_kq_k\right)+\sum_k(M-s_k)q_k
\]
forces $q_m=1$.
Since every $m$-subset has positive $\pi_m$-mass, every such subset lies in
$\mathcal F_y$. Given $i\ne j$, choose an $(m-1)$-subset $C$ disjoint from
$\{i,j\}$. Comparing $C\cup\{i\}$ and $C\cup\{j\}$ gives $h_i=h_j$.
\end{proof}


\begin{samepage}
\begin{proposition}[An atom and third-moment inequality]\label{prop:signed-atom-hermite}
Let $B_1,\ldots,B_n$ be independent Bernoulli variables with
$p_i=\Pp(B_i=1)\in(0,1)$. Let
\[
 X_i=h_i(B_i-p_i),\qquad h_i>0,\qquad
 a_i=p_i(1-p_i),\qquad t_i=1-2p_i,
\]
and put
\[
 S=\sum_iX_i,\qquad V=\sum_i a_ih_i^2,\qquad
 \kappa=\sum_i a_ih_i^3t_i,\qquad
 \rho=\frac12\sum_i a_ih_i^3(1+t_i^2).
\]
Choose $\varepsilon\in\{-1,1\}$ so that
$\varepsilon\kappa=|\kappa|$; when $\kappa\ne0$, this is
$\varepsilon=\sgn\kappa$. Set
\[
 g_i=1+\ke(t_i-\varepsilon t_*)^2,\qquad
 A=\sum_i\frac{a_i}{g_i^2}.
\]
For $\nu\geq1$ and $c\in(0,1]$, define
\[
 \mathcal M(\nu,c)
 :=\sqrt{\frac{\pi\nu c}2}\,
   \frac1\pi\int_0^\pi
   \bigl(1-c\sin^2(\theta/2)\bigr)^{\nu/2}\,d\theta .
\]
The scalar inequality
\begin{equation}
 \mathcal M(\nu,1-t^2)\leq1+\ke(t-t_*)^2
 \qquad(\nu\geq1,\ 0\leq t<1),                        \label{eq:signed-atom-scalar}
\end{equation}
is proved in Appendix~\ref{app:scalar}. It implies the joint atom and third-moment bound
\begin{equation}
 \sqrt{2\pi V}\,\frac12\sup_x\Pp(S=x)
 +\frac{|\kappa|}{6V}
 \leq\ke\,\frac{\rho}{V},                             \label{eq:atom-hermite-resource}
\end{equation}
for every finite array of centered two-point laws.
\end{proposition}
\end{samepage}

\begin{proof}
The identities $\ke t_*=1/6$ and $\ke(1-t_*^2)=1$ give
\[
 \ke\rho_i
 =a_ih_i^3\left(\frac12+\frac{\varepsilon t_i}{6}\right)
  +\frac{\ke a_ih_i^3}{2}(t_i-\varepsilon t_*)^2,
 \qquad \rho_i=\frac12a_ih_i^3(1+t_i^2).
\]
Since $\varepsilon\kappa=|\kappa|$, the
inequality~\eqref{eq:atom-hermite-resource} is equivalent to
\begin{equation}
 \sup_x\Pp(S=x)\sqrt{2\pi V}
 \leq\frac1V\sum_i a_ig_ih_i^3.                       \label{eq:physical-atom-resource}
\end{equation}
Let $q_{\max}=\max_k\Pp(N=k)$, where $N=\sum_iB_i$.
Translation and Proposition~\ref{prop:weighted-antichain} give
$\sup_x\Pp(S=x)\leq q_{\max}$. H\"older's inequality gives
\[
 V=\sum_i(a_ig_ih_i^3)^{2/3}(a_i/g_i^2)^{1/3}
 \leq\left(\sum_i a_ig_ih_i^3\right)^{2/3}A^{1/3}.
\]
Thus~\eqref{eq:physical-atom-resource} follows if
\begin{equation}
 q_{\max}\sqrt{2\pi A}\leq1.                           \label{eq:modal-atom-resource}
\end{equation}
Equality in this application of H\"older's inequality holds when $h_i$ is
proportional to $g_i^{-1}$.

Put $c_i=4a_i=1-t_i^2$, $w_i=a_i/(Ag_i^2)$, and $\nu_i=w_i^{-1}$.
Then $\sum_iw_i=1$ and $\nu_ic_i=4Ag_i^2$. Fourier inversion and
generalized H\"older give
\[
 q_{\max}\sqrt{2\pi A}
 \leq\prod_i
 \left\{\frac{\mathcal M(\nu_i,c_i)}{g_i}\right\}^{w_i}.
\]
If $\varepsilon t_i\geq0$, the scalar condition~\eqref{eq:signed-atom-scalar}
applied to $t=\varepsilon t_i$ bounds the corresponding factor by one.
If $\varepsilon t_i<0$, replace it by $|\varepsilon t_i|$; the left side is
unchanged and
\[
 1+\ke(t_i-\varepsilon t_*)^2
 =1+\ke(|t_i|+t_*)^2
 \geq1+\ke(|t_i|-t_*)^2.
\]
Hence every factor is at most one, proving the modal
bound~\eqref{eq:modal-atom-resource}.
\end{proof}


For any atom $y$, put $m_y=\Pp(S=y)$ and
$F_S^0(y)=\{F_S^+(y)+F_S^-(y)\}/2$. Then
\[
 \max\{F_S^+(y)-\Phi(y/B),\Phi(y/B)-F_S^-(y)\}
 =\frac{m_y}{2}+|F_S^0(y)-\Phi(y/B)|.
\]
Dividing \eqref{eq:atom-hermite-resource} by $\sqrt{2\pi V}$ gives
\[
 \frac12\sup_y\Pp(S=y)+\frac{|\kappa|}{6\sqrt{2\pi}V^{3/2}}
 \leq\be\ell.
\]
The second term is the uniform amplitude of
$\kappa(1-z^2)\phi(z)/(6V^{3/2})$, since
$\sup_z|(1-z^2)\phi(z)|=\phi(0)$.
A bound for the remaining midpoint CDF error at the same threshold is still
needed for the bound with arbitrary diameters in Problem~\ref{prob:two-point}.

\begin{samepage}
\begin{theorem}[Closed intervals]\label{thm:interval}
Let $Y=\sum_{i=1}^n h_iB_i$, where the $B_i$ are independent,
$0<p_i<1$, and $h_i>0$. For $r\geq0$, put
$Q(Y,r)=\sup_x\Pp\{Y\in[x,x+r]\}$.
Choose a nonempty $J\subseteq\{1,\ldots,n\}$, put $d=|J|$,
$\delta_J=\min_{i\in J}h_i$, and let $s^J_{(1)}\geq\cdots\geq s^J_{(d+1)}$
be the decreasing arrangement of the rank probabilities of
$N_J=\sum_{i\in J}B_i$. Then
\begin{equation}\label{eq:interval}
 Q(Y,r)\leq\sum_{a=1}^L s^J_{(a)},\qquad
 L=\min\{\lfloor r/\delta_J\rfloor+1,d+1\}.
\end{equation}
If $r<\delta_J$, $N_J$ has a unique mode $m\in\{1,\ldots,d-1\}$,
and equality holds, all selected $m$-subset sums lie in one interval of
length $r$. Consequently $|h_i-h_j|\leq r$ for every $i,j\in J$.
\end{theorem}
\end{samepage}
\begin{proof}
Condition on the variables outside $J$. For each outside configuration the
interval is translated, but its length remains $r$. Use the maximal chain
$A_0\subset\cdots\subset A_d$ from
Proposition~\ref{prop:weighted-antichain} on $J$, and set
$H_k=\sum_{i\in A_k}h_i$. Along each chain, $H_{k+1}-H_k\geq\delta_J$,
so a closed interval of length $r$ contains at most $L$ chain values.
If $q_k$ is the probability that $H_k$ lies in the translated interval,
then $0\leq q_k\leq1$ and $\sum_kq_k\leq L$.
The conditional interval probability is $\sum_k s_k^Jq_k$.
Choose a set $T$ of $L$ largest ranks and put $\lambda=\min_{k\in T}s_k^J$.
The exact identity
\begin{align*}
 \sum_{k\in T}s_k^J-\sum_k s_k^Jq_k
 ={}&\sum_{k\in T}(s_k^J-\lambda)(1-q_k)
 +\sum_{k\notin T}(\lambda-s_k^J)q_k
 +\lambda\left(L-\sum_kq_k\right)
\end{align*}
has nonnegative terms, also when the cutoff rank mass is tied.
Averaging over the outside configurations proves \eqref{eq:interval}.

For equality with $L=1$, every positive-probability outside configuration
must itself attain the same bound. Fix one such configuration. Uniqueness
of the mode in the preceding identity forces $q_m=1$ and $q_k=0$ for
$k\ne m$. Every $m$-subset has positive conditional rank mass, so all of its
subset sums lie in this one translated interval. For distinct $i,j\in J$,
choose an $(m-1)$-subset $C$ disjoint from them. The sums for $C\cup\{i\}$
and $C\cup\{j\}$ differ by $h_i-h_j$, proving the assertion.
\end{proof}
The strict cutoff $r<\delta_J$ is necessary for closed intervals: a single
Bernoulli variable with gap $r$ has its entire mass in $[0,r]$.
Uniqueness and interiority of the mode enter precisely in the equality
argument. A boundary mode can attain the atom bound for unequal gaps.

\begin{samepage}
\begin{corollary}[Interval concentration and selected moments]\label{cor:interval-moments}
For $S=\sum_i h_i(B_i-p_i)$, let $J\ne\varnothing$ satisfy $h_i>r$ on $J$.
Write
\[
 V_J=\sum_{i\in J}p_iq_ih_i^2,\qquad
 \kappa_J=\sum_{i\in J}p_iq_i(q_i-p_i)h_i^3,\qquad
 R_J=\sum_{i\in J}p_iq_i(p_i^2+q_i^2)h_i^3.
\]
Then
\begin{equation}\label{eq:interval-moment}
 \frac12\sqrt{2\pi V_J}\,Q(S,r)+\frac{|\kappa_J|}{6V_J}
 \leq\ke\frac{R_J}{V_J}.
\end{equation}
With $\vartheta_J=V_J/B^2$, the full-variance form is
\begin{equation}\label{eq:interval-full}
 \vartheta_J^{3/2}\frac{\sqrt{2\pi}}2Q(S,r)
 +\frac{|\kappa_J|}{6B^3}\leq\ke\frac{R_J}{B^3}.
\end{equation}
\end{corollary}
\end{samepage}
\begin{proof}
Theorem~\ref{thm:interval} gives $Q(S,r)\leq q_J:=\max_k\Pp(N_J=k)$.
The proof of \eqref{eq:atom-hermite-resource} establishes the stronger
modal inequality with $q_J$ in place of the atom mass of the selected
weighted sum. Substitution proves \eqref{eq:interval-moment}; multiplication
by $V_J/B^3$ proves \eqref{eq:interval-full}.
\end{proof}
Near equality in the rank bound~\eqref{eq:interval} with $L=1$ gives no
uniform closeness of the diameters. For example, let $0<\eta<1/2$,
$p_1=\eta$, $p_2=1/2$, $h_1=M>1+r$, $h_2=1$, and $0\leq r<1$.
The unique modal rank is $1$, of mass $1/2$, while
$Q(Y,r)=\Pp\{Y\in[1,1+r]\}=(1-\eta)/2$.
Thus the difference from the rank bound is $\eta/2$, although $M-1$ can
be arbitrarily large.

\section{Moment optimization and three-point reduction}\label{sec:moments}
At a fixed threshold, convolution is affine in each summand law. Fixing its
mean and variance therefore gives a moment optimization problem with three
affine constraints. Extreme-point theorems for moment sets, including
Winkler's characterization \cite[Theorem~2.1]{Winkler1988}, provide the broader setting.
Strong duality under interior moment constraints is also classical
\cite[Theorem~2.2]{BertsimasPopescu2005}. We give the compactness and
separation arguments for the precise coercive objectives needed here.
\begin{samepage}
\begin{theorem}[Strong quadratic moment duality]\label{thm:strongdual}
Let $s>0$ and let $h:\R\to\R$ be upper semicontinuous with
\[
 -K_-(1+|u|^3)\leq h(u)\leq K_+-\lambda|u|^3
\]
for some finite $K_-,K_+$ and $\lambda>0$. Then
\begin{equation}
 \max_{\substack{P:\ \int dP=1,\ \int u\,dP=0,\ \int u^2\,dP=s\\
                         \int|u|^3dP<\infty}}
       \int h\,dP
 =\min_{\substack{q(u)=A+Bu+Cu^2\\q\geq h\ \text{on }\R}}(A+Cs).
 \label{eq:coordinate-duality}
\end{equation}
Both extrema are attained. An extreme primal optimizer has at most three atoms, and every primal
optimizer is concentrated on the contact set of every optimal quadratic majorant.
\end{theorem}
\end{samepage}

\begin{proof}
Fix any finite-support feasible law $P_0$. A maximizing sequence may be taken with value at least
$\int h\,dP_0-1$; together with $h\leq K_+-\lambda|u|^3$, this gives uniformly bounded third moments.
The laws are tight, and the uniform bound on third moments makes
$u^2$ uniformly integrable. A weak limit preserves the three moment
constraints and has finite third moment by Fatou's lemma. The bounded-above
upper semicontinuous version of Portmanteau's theorem shows that the limit
attains the maximum. The maximizing set is weakly compact and convex; it is
a face of the feasible set because the objective is affine. It has an
extreme point, which is also extreme in the feasible set.

If an extreme feasible law $P$ had four disjoint Borel sets
$A_1,\ldots,A_4$ of positive mass, their vectors
\[
 \left(P(A_j),\int_{A_j}u\,dP,\int_{A_j}u^2\,dP\right)\in\R^3
\]
would be linearly dependent. A nonzero bounded function
$g=\sum_jc_j\ind_{A_j}$ would then satisfy
$\int g\,dP=\int ug\,dP=\int u^2g\,dP=0$.
For sufficiently small $\epsilon>0$, the two distinct laws
$(1\pm\epsilon g)P$ remain feasible and have finite third moments, while
their average is $P$. Thus an extreme law has at most three support points.

For strong duality, let $\mathcal D\subset\R^4$ consist of
\[
 \left(\mu(\R),\int u\,d\mu,\int u^2\,d\mu,z\right),
 \qquad \mu\geq0,\qquad z\leq\int h\,d\mu,
\]
where $\mu$ has finite third moment. This is a convex cone. It is closed: along a convergent
sequence with $z$ bounded below, the inequalities
$z\leq\int h\,d\mu\leq K_+\mu(\R)-\lambda\int|u|^3d\mu$ give a uniform third-moment bound; weak
compactness, uniform integrability of $u^2$, and Portmanteau for the bounded-above upper semicontinuous function
$h$ then retain all four inequalities in the limit.

Let $M$ denote the primal maximum. The point $p=(1,0,s,M)$ belongs to $\mathcal D$, whereas
$p+(0,0,0,\eps)$ does not for any $\eps>0$, so $p$ is a boundary point. A nonzero supporting
functional $(a,b,c,d)$ exists; since $\mathcal D$ is a cone and contains both $0$ and $2p$, its
supporting level at $p$ is zero. Thus
$a m_0+b m_1+c m_2+d z\leq0$ on $\mathcal D$ and equality holds at $p$. The downward vertical ray in
$\mathcal D$ gives $d\geq0$. In fact $d>0$: if $d=0$, then $(a,b,c)$ would support the moment cone
\[
 \{0\}\cup\{(m_0,m_1,m_2):m_0>0,\ m_2\geq m_1^2/m_0\}
\]
at $(1,0,s)$, which is an interior point because $s>0$, forcing the whole functional to vanish.
Normalize $d=1$. Applying the support inequality to a point mass at $u$ gives
$a+bu+cu^2+h(u)\leq0$, so $q(u):=-a-bu-cu^2$ majorizes $h$. Equality at $p$ gives
$M=-a-cs=A+Cs$. For every feasible $P$ and every feasible $q$, integration gives $\int h\,dP\leq A+Cs$. Thus this attained dual value is minimal. Finally, for any primal and dual optimizers,
$0=(A+Cs)-\int h\,dP=\int(q-h)\,dP$ with $q-h\geq0$, proving the contact assertion.
\end{proof}


\begin{samepage}
\begin{theorem}[Three-point reduction]\label{thm:threepoint}
Every strict failure of $\Delta\leq\be\ell$ can be replaced by a strict
failure with the same number of independent centered summands, each
supported on at most three points.
\end{theorem}
\end{samepage}
\begin{proof}
Fix a strict one-sided failure, its threshold $y$, its choice
$\varepsilon\in\{+1,-1\}$, and the variances of all summands.
Put $\lambda=\be B^{-3}>0$. When all coordinates except $X_i$ are fixed,
let $T_i=\sum_{j\ne i}X_j$ and set
\[
 h_i(u)=\varepsilon F_{T_i}^\varepsilon(y-u)-\lambda|u|^3.
\]
The first term is upper semicontinuous for both signs. Indeed the closed CDF
composed with $y-u$ is upper semicontinuous, while the strict CDF composed
with $y-u$ is lower semicontinuous. The cubic term gives the two bounds
required in Theorem~\ref{thm:strongdual}. A coordinate of zero variance
is deterministic. Every other coordinate can be replaced by an extreme
maximizer with its mean and variance unchanged. This increases or preserves
\[
 \varepsilon F_{\sum_iX_i}^\varepsilon(y)-\varepsilon\Phi(y/B)
 -\lambda\sum_i\E|X_i|^3.
\]
Sequential replacement gives the claimed finite-support failure.
\end{proof}

\section{The remaining finite problems}\label{sec:finite-reduction}
We formulate two open inequalities whose simultaneous proof would imply
the unrestricted sharp bound. The first is the bound for Bernoulli arrays
with arbitrary positive diameters. The second concerns the centered
three-point laws supplied by Theorem~\ref{thm:threepoint}.

\subsection{A minimal counterexample in the two-point class}
For $n\geq1$, $p\in(0,1)^n$, and $h\in(0,\infty)^n$, put
\[
 a_i=p_iq_i,\quad r_i=p_i^2+q_i^2,\quad
 \sum_i a_i h_i^2=1,\quad R=\sum_i a_ir_ih_i^3.
\]
For $A\in\{0,1\}^n$, set
\begin{equation}\label{eq:finite-deficit}
 x_A=\sum_i(A_i-p_i)h_i,\qquad
 D_A(p,h)=\be R+\Phi(x_A)
 -\sum_{C\in\{0,1\}^n}\pi_C(p)\ind_{\{(C-A)\cdot h\leq0\}},
\end{equation}
where $\pi_C(p)=\prod_i p_i^{C_i}q_i^{1-C_i}$.

\begin{samepage}
\begin{problem}[Arbitrary Bernoulli diameters]\label{prob:two-point}
Prove $D_A(p,h)\geq0$ on the complete domain in
\eqref{eq:finite-deficit}, for every $n$ and $A$.
\end{problem}
\end{samepage}
Equation~\eqref{eq:one-sided}, reflection, and scaling show that this problem
is equivalent to the sharp bound for all finite arrays of two-point laws.
The following attainment statement removes the zero-probability and
infinite-diameter boundaries from a least-dimensional failure.

\begin{samepage}
\begin{theorem}[Least-dimensional two-point failure]\label{thm:attainment}
If Problem~\ref{prob:two-point} fails, let $n$ be its least failing dimension
and let $d_n<0$ be the infimum of \eqref{eq:finite-deficit} in that dimension.
Then $d_n$ is attained at probabilities in $(0,1)$ and finite positive
diameters, with every variance share $p_iq_i h_i^2>0$.
One has $n\geq3$, the diameters are not all equal, and the selected threshold
is strictly between the two support endpoints. Subset-sum collisions may
occur.
\end{theorem}
\end{samepage}
\begin{proof}
Since $D_A\geq\be R-1$, the infimum is finite and every negative minimizing
sequence has uniformly bounded $R$. Put $w_i=a_i h_i^2$, so
$\sum_iw_i=1$. The contribution to $R$ is
$r_iw_i^{3/2}/\sqrt{a_i}$, and $r_i\geq1/2$ gives
\[
 w_i^3/a_i\leq4R^2.
\]
Pass to a subsequence on which $A$ is fixed, $p_i\to p_i^*$, and
$w_i\to w_i^*$. The set $L=\{i:w_i^*>0\}$ is nonempty.
For $i\in L$, the last inequality gives $0<p_i^*<1$, and
$h_i\to\sqrt{w_i^*/(p_i^*q_i^*)}\in(0,\infty)$.
The centered sum of coordinates outside $L$ tends to zero in $L^2$.
Thus the full sums converge weakly to the limiting sum $S_L$, which still
has variance one.

Pass to a further subsequence for which $x_A$ converges in the extended
real line. If its limit is $+\infty$ or $-\infty$, tightness and the normal
tail give a nonnegative lower limit for the deficit, a contradiction.
If $x_A\to z\in\R$, moving-threshold Portmanteau gives
\[
 \limsup\Pp(S\leq x_A)\leq\Pp(S_L\leq z).
\]
The moments of the retained coordinates converge; those of the discarded
coordinates are nonnegative. Therefore
\[
 d_n\geq\be R_L+\Phi(z)-\Pp(S_L\leq z).
\]
The right side is negative. Its CDF is then nonzero, so moving $z$ left to
the greatest retained support point below it can only decrease that deficit.
Minimality of $n$ forces $L=\{1,\ldots,n\}$.

At this interior limit, each term
$\pi_C(p)\ind_{\{(C-A)\cdot h\leq0\}}$ is upper semicontinuous.
The other terms in \eqref{eq:finite-deficit} are continuous. Consequently
$D_A$ is lower semicontinuous and its limit value is at most $d_n$.
By the definition of the infimum it equals $d_n$.
Theorem~\ref{thm:two-summands} gives $n\geq3$;
Theorem~\ref{thm:equal-gap-two-point} excludes equal diameters; and
Proposition~\ref{cor:all-support-owners} excludes support endpoints.
\end{proof}

At the attained point, variations that preserve all subset-sum collisions
satisfy the following necessary conditions. Let $M$ have independent rows spanning every subset-sum equality
active at $h$, so $Mh=0$. All other subset orders are strict in a relative
neighborhood of $h$. Set $d_i=A_i-p_i$, $t_i=1-2p_i$, $x=x_A$, and
$y=\sum_i A_i h_i$. There are multipliers $\lambda$ and $\gamma$ such that
\begin{align}
 3\be a_ir_ih_i^2+\phi(x)d_i-2\lambda a_ih_i
 +(M^{\mathsf T}\gamma)_i&=0,\label{eq:gap-stationarity}\\
 \be t_i^3h_i^3-\phi(x)h_i+I_i-\lambda t_i h_i^2&=0,
 \qquad I_i=\Pp\{y-h_i<Y_{-i}\leq y\},\label{eq:prob-stationarity}\\
 2\lambda&=3\be R+x\phi(x),\label{eq:multiplier}
\end{align}
where $Y_{-i}=\sum_{j\ne i}h_jB_j$. Moreover, for every $v$ with
$Mv=0$ and $\sum_i a_ih_iv_i=0$,
\begin{equation}\label{eq:gap-hessian}
 \sum_i(6\be a_ir_ih_i-2\lambda a_i)v_i^2
 -x\phi(x)\left(\sum_i d_iv_i\right)^2\geq0.
\end{equation}
To verify these conditions, apply the Lagrange multiplier rule to
$D_A-\lambda(\sum_i a_ih_i^2-1)+\gamma\cdot Mh$ on the fixed stratum.
The constraint qualification holds because $Mh=0$, whereas the variance
gradient has scalar product $2$ with $h$. Differentiation uses
$(a_ir_i)'=t_i^3$, $a_i'=t_i$, and $\partial_{p_i}F_S^+(x_A)=-I_i$.
Multiplication of \eqref{eq:gap-stationarity} by $h_i$ and summation gives
\eqref{eq:multiplier}; the gap Hessian is \eqref{eq:gap-hessian}.
Changing the representative $A$ of a colliding atom adds a vector in the
row span of $M$. Its contribution is absorbed by $\gamma$, and its scalar
product with every allowed $v$ is zero.
A proof that no negative point satisfies this system on any complete
collision stratum would solve Problem~\ref{prob:two-point}.

\subsection{An open inequality for three-point mixtures}
For $a,b>0$, define the centered pair
\[
 Q_{a,b}=\frac b{a+b}\delta_{-a}+\frac a{a+b}\delta_b,
 \quad \Var(Q_{a,b})=ab,\quad
 \int|u|^3\,dQ_{a,b}(u)=\frac{ab(a^2+b^2)}{a+b}.
\]
A centered law with exactly three support points has one of the
representations
\begin{equation}\label{eq:canonical-pairs}
 P=\alpha\delta_0+(1-\alpha)Q_{a,b},
 \qquad\hbox{or, up to reflection,}\qquad
 P=\alpha Q_{a,b_1}+(1-\alpha)Q_{a,b_2},\quad 0<b_1<b_2,
\end{equation}
with $0<\alpha<1$. The first formula follows from centering on
$\{-a,0,b\}$. For masses $m_-,m_1,m_2$ at $-a,b_1,b_2$, centering gives
$am_-=b_1m_1+b_2m_2$; the mixture weights are
$m_j(a+b_j)/a$ and sum to one.
If $P$ is an extreme optimizer in Theorem~\ref{thm:strongdual}, every
point in \eqref{eq:canonical-pairs} is a contact of the same optimal
quadratic. Each component law therefore maximizes the same coordinate objective
among centered laws of its own variance. Indeed, at variance $r$ the
quadratic bounds the objective by $A+Cr$, and a component supported on
contacts attains this value.

Fix $T=\sum_jT_j$, where the $T_j$ are independent centered variables
with at most three atoms each, and put $V_T=\Var(T)>0$ and
$R_T=\sum_j\E|T_j|^3$.
Choose $\varepsilon\in\{+1,-1\}$ and $y\in\R$. Let
$P=\alpha Q_-+(1-\alpha)Q_+$ be one of \eqref{eq:canonical-pairs}, and write
\[
 r_-=\Var(Q_-)<r_+=\Var(Q_+),\quad
 s=\alpha r_-+(1-\alpha)r_+,\quad V=V_T+s,\quad\lambda=\be V^{-3/2}.
\]
Suppose $P$ is an extreme maximizing law, of variance $s$, for
$h(u)=\varepsilon F_T^\varepsilon(y-u)-\lambda|u|^3$.
For independent $X_\pm\sim Q_\pm$, define
\[
 V_\pm=V_T+r_\pm,\quad R_\pm=R_T+\E|X_\pm|^3,\quad
 R=\alpha R_-+(1-\alpha)R_+,
\]
\begin{equation}\label{eq:endpoint-delta}
 \delta_\pm=\varepsilon\Phi(y/\sqrt{V_\pm})
 +\be R_\pm V_\pm^{-3/2}
 -\varepsilon F_{T+X_\pm}^\varepsilon(y).
\end{equation}
With $X\sim P$ independent of $T$, the exact excess of the mixture is
\begin{equation}\label{eq:mixing-identity}
 \varepsilon F_{T+X}^\varepsilon(y)-\varepsilon\Phi(y/\sqrt V)
 -\be RV^{-3/2}=G-\alpha\delta_--(1-\alpha)\delta_+,
\end{equation}
where
\begin{align}\label{eq:mixing-G}
 G={}&\varepsilon\left\{\alpha\Phi(y/\sqrt{V_-})
 +(1-\alpha)\Phi(y/\sqrt{V_+})-\Phi(y/\sqrt V)\right\}\notag\\
 &+\be\left\{\alpha R_-V_-^{-3/2}+(1-\alpha)R_+V_+^{-3/2}
 -RV^{-3/2}\right\}.
\end{align}
This follows by adding and subtracting the endpoint normal approximations;
both convolution and the third absolute moment are affine in the mixture.

\begin{samepage}
\begin{problem}[Three-point mixture]\label{prob:mixture}
In the complete domain just specified, suppose that both endpoint arrays
$T+X_-$ and $T+X_+$ satisfy the sharp Berry--Esseen bound at their respective
variances. Prove
\begin{equation}\label{eq:mixing-problem}
 \alpha\delta_-+(1-\alpha)\delta_+\geq G.
\end{equation}
\end{problem}
\end{samepage}

\begin{samepage}
\begin{proposition}[Reduction to the two finite problems]\label{prop:final-reduction}
A solution of Problems~\ref{prob:two-point} and~\ref{prob:mixture} implies
$\Delta\leq\be\ell$ for every independent array in
\eqref{eq:normalization}. If the two-point problem is solved but the general
bound fails, Problem~\ref{prob:mixture} has a counterexample with $G>0$.
\end{proposition}
\end{samepage}
\begin{proof}
Use Theorem~\ref{thm:threepoint} and choose a strict counterexample with the
fewest coordinates having exactly three atoms. Under
Problem~\ref{prob:two-point}, this number is positive.
The one-summand estimate of Bentkus and Kirsa
\cite[Theorem~1]{BentkusKirsa1989} gives
$\Delta\leq c_{\mathrm{BK}}\ell$ with
$c_{\mathrm{BK}}=0.370352\ldots<\be$ for every centered law with finite
absolute third moment. It excludes a single nondegenerate coordinate; hence deleting any chosen
three-point coordinate leaves positive variance.
Optimize that coordinate at its original variance and the selected sign and
threshold. An extreme optimizer exists by Theorem~\ref{thm:strongdual}.
It cannot have two atoms, since that would give a strict failure with fewer
three-point coordinates. Thus it has exactly three atoms and admits
\eqref{eq:canonical-pairs}. Each endpoint replacement removes one three-point
coordinate. By minimality, each endpoint array satisfies the sharp bound,
so $\delta_-,\delta_+\geq0$ at the selected threshold. The optimized mixture
still has positive excess. Equation~\eqref{eq:mixing-identity} now gives
$G>\alpha\delta_-+(1-\alpha)\delta_+\geq0$, contradicting
\eqref{eq:mixing-problem}.
\end{proof}

\section{Computational proof supplement}\label{sec:verification}
Five primary calculations are used in the proofs: the exceptional-Bernoulli
curvature inequality \eqref{eq:one-exceptional-compact-gap}, the two
inequalities of Lemma~\ref{lem:two-compact}, the rational odds comparison
in Appendix~\ref{app:support}, and the scalar Fourier comparison in
Appendix~\ref{app:scalar}. Their mathematical domains and analytic complements
are part of the proofs. The programs under \path{verify/finite/} implement
the finite parts. The reported primary results and source hashes are in
\path{certificates/finite/results.json}; \path{certificates/finite/README.md}
lists the programs and rerun commands. Separate evaluators are supplied for
the scalar atom inequality, the odds comparison, and the first interior
two-summand threshold. The second threshold has a higher-precision rerun
and separate formula and boundary checks. Paths are relative to the
public supplement root.

The interval calculations use Arb through \texttt{python-flint} 0.8.0.
Arb encloses real quantities by balls with certified outward rounding,
including the elementary and special functions used here
\cite{Johansson2017}. A numerical estimate without an enclosure is never an
acceptance test. Every rational input box is enclosed; a lower bound is
accepted only if its complete outward interval is positive. The use of
floating-point scores to choose a subdivision direction in one program
cannot change this test or the coverage identity.

\subsection{The two-summand covers}
For \eqref{eq:two-upper-middle},
\path{verify/finite/two_summands/certificate.py} starts from $[0,1]^3$
in $(p_1,p_2,t)$ and uses 192-bit arithmetic. For a probability interval
$[l,u]$, its variance range has lower endpoint
$\min\{l(1-l),u(1-u)\}$ and upper endpoint $1/4$ if the interval contains
$1/2$, otherwise $\max\{l(1-l),u(1-u)\}$.
The program obtains a lower bound for $R/V^{3/2}$ by taking the greater of
direct interval evaluation, when $V$ is bounded away from zero, and
\eqref{eq:reciprocal-gap-resource}. In the latter bound it uses the upper
enclosure of $a_1/r_1^2+a_2/r_2^2$ before taking its reciprocal square root.
It subtracts an upper enclosure of $1-p_1p_2$.

For the normal term, direct evaluation applies when the variance interval
has positive lower endpoint. If the variance interval meets zero and the
numerator has nonnegative lower endpoint, dividing that endpoint by the
upper variance square root gives a valid lower bound for the argument.
A sign-indefinite numerator receives the lower CDF bound zero.
The additional inequality
\[
 \frac{-p_1t+q_2}{\sqrt{a_1t^2+a_2}}\geq-\sqrt{p_1/q_1}
\]
provides a second bound. If the numerator is negative, its magnitude is
at most $p_1t$ and the denominator is at least $\sqrt{a_1}t$; otherwise
the claim is immediate. The analogous bound
$(q_1t-p_2)/\sqrt V\geq-\sqrt{p_2/q_2}$ is used by the accompanying
checks. On boxes touching the corresponding deterministic probability,
the limiting CDF lower bound zero avoids a singular odds evaluation.
Thus every bound remains valid on the positive-variance part of a box
which touches a boundary face.

A box failing the strict positivity test is bisected into two closed
rational boxes along one coordinate. The adaptive choice of coordinate
uses lower bounds only as a heuristic; both children are retained.
The complete calculation processed $2231$ boxes and accepted $1116$,
with maximum accepted depth $16$. The smallest accepted outward lower
endpoint exceeded $3.2325447770\times10^{-5}$.

For \eqref{eq:two-lower-middle},
\path{verify/finite/two_summands/row10_certificate.py} starts from
$[0,1]^3$ in $(p_1,p_2,2\theta/\pi)$ and uses 192-bit arithmetic.
The shape
\[
 s(p)=\frac{p^2+(1-p)^2}{\sqrt{p(1-p)}}
\]
is decreasing on $(0,1/2]$ and increasing on $[1/2,1)$.
For a probability interval its infimum is therefore $1$ if the interval
contains $1/2$, and otherwise its value at the endpoint nearest $1/2$.
The two nonnegative third-moment terms are bounded separately using lower
bounds for $\cos\theta$, $\sin\theta$, and $s(p_i)$.
For the argument of $\Phi$, the positive term uses the lower bound for
$\cos\theta$ and the largest $p_1$; the negative term uses the upper bound
for $\sin\theta$ and the largest $p_2$. If the latter endpoint is $1$,
the CDF lower bound is zero. If the former endpoint is $1$, its positive
term is replaced by zero. This treats all boxes meeting probability
boundaries without assigning values to an undefined standardized law.
The upper bound for the subtracted $q_2$ uses the smallest $p_2$.
The wider coordinate is bisected whenever the bound is inconclusive.
The complete cover processed $3509$ boxes and accepted $1755$, with
maximum depth $17$ and smallest outward lower endpoint greater than
$1.3393433322\times10^{-5}$.

Each run starts with one box and replaces every unaccepted box by both
children. Its terminal queue is empty, and its counts satisfy
$N_{\rm processed}=2N_{\rm accepted}-1$.
Depth or box-budget exhaustion terminates with an error rather than a
positive conclusion. These facts establish coverage of the whole cubes.
The first algebraic inequality extends to every positive-variance cube
point. Its interpretation as the middle closed CDF requires $t>0$.
At $t=0$ the correct distribution is obtained by deleting the zero-diameter
coordinate. Probability endpoints are likewise handled by deletion in the
proof of Theorem~\ref{thm:two-summands}.

\subsection{The other finite calculations}
The exceptional-Bernoulli calculation uses three four-dimensional unit-box
charts displayed in the proof of
Lemma~\ref{lem:one-exceptional-compact}. It starts with one box per chart;
the recorded complete cover has $448835$ processed boxes and $224419$
accepted boxes, satisfying $N_{\rm processed}=2N_{\rm accepted}-3$.
Its maximum accepted depth is $31$, and the smallest accepted outward
lower endpoint exceeds $1.8532764256\times10^{-7}$.
The separate analytic inequalities in that proof cover small total
variance, binomial variance $W\geq3$, and nonpositive adjacent mass
differences. The direct two-summand proof above supplies the remaining
case $M=1$.

The scalar atom calculation is a fixed cover of $63$ order slabs by
$2048$ intervals each. Formula~\eqref{eq:scalar-slab-majorant} bounds every
real order in each slab; a grid of real orders is not used. The positive
polynomial \eqref{eq:scalar-even-moment} is evaluated in Arb, and all
$129024$ boxes satisfy the strict bound
\eqref{eq:scalar-compact-gap}. An independently written evaluator uses
direct coefficient sums instead of Horner evaluation. The infinite-order
complement is reduced analytically to the single Bessel inequality stated
in Appendix~\ref{app:scalar} and included in the verification material.

The odds comparison uses SymPy \cite{Meurer2017} for exact arithmetic in
$\mathbb Q(\sqrt5)$.
For a polynomial $P(r,t)=\sum c_{ij}r^it^j$ and degrees $D,E$, its tensor
Bernstein coefficients on $[0,1]^2$ are
\begin{equation}\label{eq:bernstein-conversion}
 b_{kl}=\sum_{i\leq k,\,j\leq l}
 c_{ij}\frac{\binom{k}{i}}{\binom{D}{i}}
       \frac{\binom{l}{j}}{\binom{E}{j}}.
\end{equation}
The identity
$P=\sum_{k,l}b_{kl}\binom{D}{k} r^k(1-r)^{D-k}\binom{E}{l} t^l(1-t)^{E-l}$
turns coefficient signs into signs on the whole square. All signs in
$\mathbb Q(\sqrt5)$ can be determined by rational arithmetic, squaring
only after the signs of the rational parts are known. Appendix~\ref{app:support}
specifies the two mapped polynomials, the degrees, and every zero coefficient.
The independent checker cancels a positive polynomial factor before a
second basis conversion.

\appendix
\section{Support endpoint estimates}\label{app:support}
Use the Bernoulli representation from Section~\ref{sec:two}. In this
appendix $V=B^2$ and $\mathcal B_h=B^3\ell$.
\begin{samepage}
\begin{lemma}[Normal tails at the support endpoints]
\label{lem:support-owners-without-atom-term}
Let $n\geq1$, and let $B_1,\ldots,B_n$ be independent, with
$B_i\sim\Bern(p_i)$, $p_i\in(0,1)$, and $h_i>0$.
Put $q_i=1-p_i$ and define
\[
 V=\sum_i p_iq_ih_i^2,
 \qquad
 \mathcal B_h=\sum_i p_iq_i(p_i^2+q_i^2)h_i^3,
\]
and set
\[
 u=\sum_iq_ih_i,
 \qquad
 m=\sum_ip_ih_i.
\]
Then
\begin{equation}
 \be\frac{\mathcal B_h}{V^{3/2}}
 >\max\left\{
  1-\Phi\left(\frac{u}{\sqrt V}\right),
  \Phi\left(-\frac{m}{\sqrt V}\right)
 \right\}.
\label{eq:support-owner-tail-bound}
\end{equation}
Consequently, the closed deficit at the upper support and the strict deficit
at the lower support are both strictly positive.
\end{lemma}
\end{samepage}

\begin{proof}
It is enough to prove the first inequality in
\eqref{eq:support-owner-tail-bound}; complementing every $B_i$ exchanges $p_i$
and $q_i$, preserves $V$ and $\mathcal B_h$, and gives the second. Normalize
$V=1$ and write
\[
 x_i=q_ih_i,\qquad v_i=p_iq_ih_i^2,\qquad
 b_i=p_iq_i(p_i^2+q_i^2)h_i^3,\qquad x=\sum_i x_i.
\]
For $c=2(\sqrt2-1)$,
\[
 b_i=\frac{p_i^2+q_i^2}{p_i}\frac{v_i^2}{x_i},
 \qquad
 \frac{p_i^2+q_i^2}{p_i}-c
 =\frac{(\sqrt2p_i-1)^2}{p_i}\geq0.
\]
Since $\sum_i v_i=1$, Cauchy--Schwarz gives
\begin{equation}
 \mathcal B_h\geq c\sum_i\frac{v_i^2}{x_i}
 \geq\frac{c}{x}.
\label{eq:support-owner-resource-bound}
\end{equation}
We claim that
\begin{equation}
 x\{1-\Phi(x)\}<\be c
 \qquad(x>0).
\label{eq:support-owner-normal-tail}
\end{equation}
If $0<x\leq2/3$, then the left side is at most $1/3$. The elementary bounds
\[
 \sqrt{10}>\frac{79}{25},\qquad
 \sqrt2>\frac{169}{120},\qquad
 \pi<\frac{22}{7},
\]
give $\be>49/120$ and $c>49/60$. For the first comparison,
\[
 20(\sqrt{10}+3)>\frac{616}{5}
 >49\sqrt{\frac{44}{7}}>49\sqrt{2\pi};
\]
the middle inequality follows on squaring, with difference $2156/25$.
The bound on $\pi$ follows, for example, from
\[
 \frac{22}{7}-\pi
 =\int_0^1\frac{t^4(1-t)^4}{1+t^2}\,dt>0.
\]
Hence $\be c>2401/7200>1/3$.
For $x\geq2/3$, Mills' inequality and monotonicity of $\ph$ give
\[
 x\{1-\Phi(x)\}<\ph(x)\leq\ph(2/3)=\ph(0)e^{-2/9}.
\]
Now $e^{-2/9}<9/11<c$, where the last inequality follows from
$\sqrt2>31/22$, and $\be/\ph(0)=(\sqrt{10}+3)/6>1$. This proves
\eqref{eq:support-owner-normal-tail}. Combining it with
\eqref{eq:support-owner-resource-bound} proves the upper-support inequality.
At that support the closed CDF equals one. At the lower support the strict CDF
equals zero, and exact complementation exchanges the two CDF limits. The two
deficit statements follow.
\end{proof}

\begin{samepage}
\begin{proposition}[An inequality for Bernoulli odds]
\label{prop:product-normalized-odds-envelope}
Let $o_1,\ldots,o_n\geq0$, not all zero, and put
\begin{equation}
 s=\sum_i o_i,\qquad
 \Pi=\prod_i(1+o_i)^{-1},\qquad
 g(o)=\frac{o(1+o)^2}{(1+o^2)^2},\qquad
 \mathcal A=\sum_i g(o_i).
\label{eq:support-odds-data}
\end{equation}
Then
\begin{equation}
 \Pi^2\mathcal A\leq\frac13\qquad(0<s\leq1),
\label{eq:support-odds-small-envelope}
\end{equation}
\begin{equation}
 \Pi^2\mathcal A\leq
 \frac{s}{(1+s/2)^2(1+s^2/4)^2}
 \qquad\left(1\leq s\leq\frac{16}{9}\right),
\label{eq:support-odds-live-envelope}
\end{equation}
and
\begin{equation}
 \Pi^2\mathcal A\leq\frac{944784}{6076225}
 \qquad\left(s\geq\frac{16}{9}\right).
\label{eq:support-odds-large-envelope}
\end{equation}
For every fixed nonzero direction $v\in[0,\infty)^n$, set
$\mathcal F_v(t)=\Pi(tv)^2\mathcal A(tv)$. Then
\[
 \frac{d\log\mathcal F_v(t)}{d\log t}<-\frac3{100}
 \qquad\left(t\sum_i v_i\geq\frac{16}{9}\right).
\]
In particular, $\mathcal F_v$ is strictly decreasing on this range.
Equality in \eqref{eq:support-odds-live-envelope} holds exactly when two
odds are positive and equal to $s/2$, up to zero coordinates and
permutation.
\end{proposition}
\end{samepage}

\begin{proof}
We first prove the intermediate-range envelope. After relabeling, isolate
$x=o_1$ and $y=o_2$, and write
\[
 a=x+y,\qquad k=xy,
 \qquad 0<k\leq\frac{a^2}{4}.
\]
For the untouched coordinates, set
\[
 H=\sum_{j=3}^n g(o_j),
 \qquad
 \Pi_0=\prod_{j=3}^n(1+o_j)^{-1}.
\]
With
\[
 z_a(k)=\frac1{(1+a+k)^2},\qquad
 R_a(k)=z_a(k)\{g(x)+g(y)\},
\]
the full functional retains the untouched coordinates exactly:
\begin{equation}
 \Pi^2\mathcal A=\Pi_0^2\{R_a(k)+Hz_a(k)\}.
\label{eq:support-odds-affine-environment}
\end{equation}
Put $D=1+a^2-2k+k^2$. Direct symmetric expansion gives
\[
 R_a(k)=\frac{N(a,k)}{D^2(1+a+k)^2},
\]
where
\[
\begin{aligned}
 N(a,k)={}&a+2a^2+a^3+(a^3-a-4)k\\
 &+(2a^2-a+8)k^2+(a-4)k^3.
\end{aligned}
\]
The pair endpoints are a merge and an equal split. Their values are
\[
 \mathcal C_1(a)=\frac{a}{(1+a^2)^2},
 \qquad
 \mathcal C_2(a)=
 \frac{a}{(1+a/2)^2(1+a^2/4)^2},
\]
and
\[
 \mathcal C_2(a)-\mathcal C_1(a)
 =-
 \frac{a^2(a^2-6a+4)(a^3+10a^2+4a+16)}
 {(a+2)^2(a^2+1)^2(a^2+4)^2}.
\]
Thus the endpoints exchange order at
$a_*=3-\sqrt5$.

Two exact Bernstein certificates control the whole pair curve. For
$0\leq a\leq a_*$, define
\[
 \mathcal S_{\mathrm{raw}}(a,k)
 :=\frac{aD^2(1+a+k)^2-(1+a^2)^2N(a,k)}{k},
 \qquad
 \mathcal S(a,k):=\frac{\mathcal S_{\mathrm{raw}}(a,k)}{(1+a)^2}.
\]
Both quotients are polynomials, and
\[
 R_a(0)-R_a(k)
 =\frac{k\mathcal S_{\mathrm{raw}}(a,k)}
 {(1+a^2)^2D^2(1+a+k)^2}.
\]
Under
\[
 a=a_*r,\qquad k=\frac{a^2t}{4},\qquad (r,t)\in[0,1]^2,
\]
the uncancelled mapped quotient has monomial bidegree $(13,5)$. Cancelling
the positive common factor $(1+a)^2$ gives a reduced quotient of bidegree
$(11,5)$. Their degree-$(60,20)$ Bernstein tables differ but have the same
sign pattern: exactly $1280$ entries are positive, and the sole zero is the
corner coefficient $(60,20)$. Hence
\[
 R_a(k)\leq R_a(0),
\]
strictly for $k>0$ except at $(a,k)=(a_*,a_*^2/4)$. Since $z_a$ decreases,
the same merge inequality holds after the term $Hz_a(k)$ in
\eqref{eq:support-odds-affine-environment} is added.

For $a_*\leq a\leq16/9$, define $\theta\in[0,1]$ by
\[
 z_a(k)=\theta z_a(0)+(1-\theta)z_a(a^2/4).
\]
Define $\mathcal L(a,k)$ by the exact identity
\[
\begin{aligned}
 &\theta R_a(0)+(1-\theta)R_a(a^2/4)-R_a(k)\\
 &\quad=
 \frac{k(a^2-4k)\mathcal L(a,k)}
 {(1+a^2)^2(4+a^2)^2(1+a+k)^2
  (8+8a+a^2)D^2}.
\end{aligned}
\]
Under
\[
 a=a_*+\left(\frac{16}{9}-a_*\right)r,
 \qquad k=\frac{a^2t}{4},
\]
the remaining quotient has bidegree $(13,4)$ and all $70$ tensor Bernstein
coefficients are positive. Therefore
\begin{equation}
 R_a(k)<\theta R_a(0)+(1-\theta)R_a(a^2/4)
 \qquad\left(0<k<\frac{a^2}{4}\right).
\label{eq:support-odds-pair-chord}
\end{equation}
The untouched term $Hz_a(k)$ is affine in the same coordinate, so
\eqref{eq:support-odds-pair-chord} holds for the complete pair profile,
uniformly in $H\geq0$.

Fix $s\in[1,16/9]$ and maximize $\Pi^2\mathcal A$ on the closed odds
simplex. No
positive pair in a maximizer can have sum at most $a_*$. For $a<a_*$,
merging the pair strictly increases the objective. At $a=a_*$, this is also
true for an equal pair when $H>0$; if $H=0$, then $s=a_*<1$. Every remaining
pair lies in the strict chord region, so it must be equal. Thus all positive
odds in any maximizer are equal. If their number is $d\geq2$, then
\[
 d<\frac{2s}{a_*}
 \leq\frac89(3+\sqrt5)<5.
\]
The case $d=1$ is already in the resulting list, so in all cases
$d\in\{1,2,3,4\}$.
The equal-block value is
\[
 \mathcal C_d(s)=
 \frac{s}{(1+(s/d)^2)^2(1+s/d)^{2(d-1)}}.
\]
After positive factors are removed, the three differences
$\mathcal C_2-\mathcal C_d$, $d=1,3,4$, have the respective factors
\[
 -(s^2-6s+4),\qquad
 8s^3-33s^2-18s+108,\qquad
 s^4+12s^3-64s^2+256.
\]
They are strictly positive on the intermediate range. The first polynomial inside
the minus sign decreases there and starts at $-1$. The second decreases there
and has right-endpoint value $12140/729$.
The third is at least $256 - 64(16/9)^2>0$. This proves
\eqref{eq:support-odds-live-envelope} and its equality statement.

For the small range, expand term by term:
\[
 \Pi^2\mathcal A=
 \sum_i\frac{o_i}
 {(1+o_i^2)^2\prod_{j\ne i}(1+o_j)^2}.
\]
If $r_i=s-o_i$, then the product over $j\ne i$ is at least $1+r_i$. The
elementary inequality
\begin{equation}
 (1+x^2)^2(1+y)^2\geq3(x+y)
 \qquad(x>0,\ y\geq0)
\label{eq:support-odds-termwise-denominator}
\end{equation}
therefore bounds every summand with $o_i>0$ by $o_i/(3s)$; a zero-odds
summand vanishes. To verify
\eqref{eq:support-odds-termwise-denominator}, put $A=(1+x^2)^2$. The
difference is
\[
 Ay^2+(2A-3)y+(A-3x).
\]
Its constant term is positive because
$\min_{x>0}A/x=16/(3\sqrt3)>3$. If $2A-3<0$, then $x<1/2$, and four times
the leading coefficient times the quadratic minimum is
$12A(1-x)-9>0$, since $A(1-x)\geq25/32$. This proves
\eqref{eq:support-odds-small-envelope}.

It remains to prove the large-range envelope. Discard the zero coordinates
of a fixed nonnegative direction. Along the resulting positive ray
$o_i=tv_i$, let $\mathcal F(t)=\Pi(tv)^2\mathcal A(tv)$ and put
\[
 e(o)=\frac{o g'(o)}{g(o)}.
\]
Logarithmic differentiation gives
\[
 \frac{d}{d\log t}\log\mathcal F(t)
 =\frac{\sum_i g(o_i)e(o_i)}{\mathcal A}
 -2\sum_i\frac{o_i}{1+o_i}.
\]
The exact identities
\[
 \frac54-e(o)
 =\frac{9o^3+17o^2-7o+1}{4(1+o)(1+o^2)}>0,
\]
\[
 17o^2-7o+1
 =17\left(o-\frac7{34}\right)^2+\frac{19}{68},
\]
and
\[
 \frac{x}{1+x}+\frac{y}{1+y}-\frac{x+y}{1+x+y}
 =\frac{xy(x+y+2)}{(1+x)(1+y)(1+x+y)}>0
\]
imply
\[
 \frac{d}{d\log t}\log\mathcal F(t)
 <\frac54-\frac{2s}{1+s}\leq-\frac3{100}
 \qquad\left(s\geq\frac{16}{9}\right).
\]
Scale back to total odds $16/9$ and apply
\eqref{eq:support-odds-live-envelope}. This gives
\eqref{eq:support-odds-large-envelope}.

The two tables are generated by
\path{verify/finite/support_odds/exact_certificate.py}; the independent
basis conversion is in the corresponding \texttt{independent\_certificate.py}.
Equation~\eqref{eq:bernstein-conversion} specifies the conversion, and
the programs check every coefficient sign. Their source hashes and the
reported sign counts are in \path{certificates/finite/results.json}.
\end{proof}


\begin{samepage}
\begin{lemma}[Endpoint atom comparison]
\label{lem:atom-bearing-support-owners}
Under the hypotheses and notation of
Lemma~\ref{lem:support-owners-without-atom-term}, the closed deficit at the
lower support and the strict deficit at the upper support are strictly
positive.
\end{lemma}
\end{samepage}

\begin{proof}
Normalize $V=1$, put $o_i=p_i/q_i$, and set
\[
 w_i=p_iq_ih_i^2,\qquad \sum_iw_i=1.
\]
The distance to the lower support, its atom, and the cubic sum are
\[
 x_0=\sum_i\sqrt{o_iw_i},\qquad
 \Pi=\prod_i(1+o_i)^{-1},
\]
\[
 \mathcal B_h
 =\sum_i\frac{1+o_i^2}{\sqrt{o_i}(1+o_i)}w_i^{3/2}.
\]
Use $s$ and $\mathcal A$ from \eqref{eq:support-odds-data}. With
\[
 a_i=\frac{o_i}{(1+o_i)^2},\qquad
 r_i=\frac{1+o_i^2}{(1+o_i)^2},\qquad
 z_i=\sqrt{a_i}h_i,
\]
we have $\sum_i z_i^2=1$ and
$\mathcal B_h=\sum_i(r_i/\sqrt{a_i})z_i^3$. Cauchy--Schwarz, twice, gives
\[
 1\leq
 \mathcal B_h\sum_i\frac{\sqrt{a_i}}{r_i}z_i
 \leq\mathcal B_h\sqrt{\mathcal A},
 \qquad
 x_0^2\leq s.
\]
Let $Q(x)=1-\Phi(x)$. Since $Q$ decreases, it is enough to prove
\begin{equation}
 \frac{\be}{\sqrt{\mathcal A}}+Q(\sqrt s)>\Pi.
\label{eq:atom-bearing-weights-free-target}
\end{equation}
Put
\begin{equation}
 f(s):=(1+s)Q(\sqrt s),\qquad f'(s)<0\quad(s>0).
\label{eq:atom-bearing-tail-monotonicity}
\end{equation}
To prove the derivative sign, take $x=\sqrt s$, $M(x)=Q(x)/\ph(x)$, and
$u(x)=(1+x^2)/(2x)$, the required derivative sign is $M(x)<u(x)$. For
$x\geq1$ this follows from the Mills bound $M(x)<1/x$. On $(0,1)$, put
$d=u-M$. We have $d(0+)=\infty$ and $d(1)>0$. Since
$M'=xM-1$, at any interior critical point of $d$,
\[
 M(x)=\frac{3x^2-1}{2x^3},
 \qquad
 d(x)=\frac{(1-x^2)^2}{2x^3}>0.
\]
Thus $d$ cannot have a nonpositive minimum, proving
\eqref{eq:atom-bearing-tail-monotonicity}.

First suppose $0<s\leq1$. By
\eqref{eq:support-odds-small-envelope},
\[
 \frac{\be}{\Pi\sqrt{\mathcal A}}>\frac7{10},
\]
using $\be>49/120$ and $\sqrt3>12/7$. Also
$\Pi\leq(1+s)^{-1}$. The alternating exponential bound on $[0,1]$ and
$\ph(0)<2/5$ give
\[
 Q(1)>\frac12-\frac25\left(1-\frac16+\frac1{40}\right)
 =\frac{47}{300}.
\]
By \eqref{eq:atom-bearing-tail-monotonicity},
\[
 \frac{Q(\sqrt s)}\Pi\geq(1+s)Q(\sqrt s)
 \geq2Q(1)>\frac{47}{150}>\frac3{10}.
\]
The two strict bounds prove
\eqref{eq:atom-bearing-weights-free-target} in the small range.

For $1\leq s\leq16/9$, write $u=s/2$. The intermediate-range envelope gives
\begin{equation}
 \frac{\be}{\Pi\sqrt{\mathcal A}}
 \geq\be\frac{(1+u)(1+u^2)}{\sqrt{2u}}
 >\frac{49}{64}.
\label{eq:atom-bearing-live-resource-payment}
\end{equation}
The factor after $\be$ increases for $u\geq1/2$; its logarithmic derivative
has the sign of $5u^3+3u^2+u-1>0$. At $s=16/9$, the alternating bound
\[
 e^{-t^2/2}\leq
 1-\frac{t^2}{2}+\frac{t^4}{8}-\frac{t^6}{48}+\frac{t^8}{384}
 \qquad\left(0\leq t\leq\frac43\right)
\]
and $\ph(0)<2/5$ give
\[
 Q\left(\frac43\right)>\frac{16727087}{186004350},
 \qquad
 \frac{25}{9}Q\left(\frac43\right)-\frac{15}{64}
 >\frac{33055039}{2142770112}>0.
\]
Consequently,
\begin{equation}
 \frac{Q(\sqrt s)}\Pi\geq(1+s)Q(\sqrt s)>\frac{15}{64}.
\label{eq:atom-bearing-live-tail-payment}
\end{equation}
The bounds in \eqref{eq:atom-bearing-live-resource-payment} and
\eqref{eq:atom-bearing-live-tail-payment} again sum to more than one.

Finally, if $s\geq16/9$, Proposition~\ref{prop:product-normalized-odds-envelope}
and the exact margin
\[
 \frac{2401}{14400}-\frac{944784}{6076225}
 =\frac{7873013}{699981120}>0
\]
give $\Pi^2\mathcal A<\be^2$. The cubic term alone proves
\eqref{eq:atom-bearing-weights-free-target}. Hence, in every range,
\[
 \be\mathcal B_h+Q(x_0)-\Pi
 \geq\frac{\be}{\sqrt{\mathcal A}}+Q(\sqrt s)-\Pi>0.
\]
This is the closed deficit at the lower support. Complementing every $B_i$
preserves $V$ and $\mathcal B_h$ and carries it to the strict deficit at the
upper support.
\end{proof}


\begin{proof}[Proof of Proposition~\ref{cor:all-support-owners}]
The two preceding endpoint lemmas supply the two signs at each endpoint.
\end{proof}

\section{The scalar Fourier inequality}\label{app:scalar}
We prove the scalar input \eqref{eq:signed-atom-scalar} used in the atom
and interval estimates. Write
\[
 I_0(W)=\frac1\pi\int_0^\pi e^{W\cos\theta}\,d\theta,
 \qquad I_1(W)=I_0'(W).
\]
For $W=\nu c/4$, the inequality $1-u\leq e^{-u}$ gives
\begin{equation}\label{eq:bessel-majorant}
 \mathcal M(\nu,c)\leq\sqrt{2\pi W}\,e^{-W}I_0(W).
\end{equation}
\begin{samepage}
\begin{proposition}[Analytic ranges of the scalar atom bound]\label{prop:scalar-atom-ranges}
The scalar inequality~\eqref{eq:signed-atom-scalar} holds in each of the
following ranges:
\[
 1\leq\nu\leq2,\quad 0\leq t<1;
 \qquad
 \nu\geq1,\quad 0\leq t\leq\frac1{\sqrt{10}};
 \qquad
 \nu\geq1,\quad t_0\leq t<1,
\]
where
\[
 t_0:=t_*+\sqrt{\frac1{5\ke}}=0.6035634\ldots.
\]
The only region not covered by these analytic estimates is
\begin{equation}
 \nu>2,\qquad \frac1{\sqrt{10}}<t<t_0.                 \label{eq:scalar-atom-strip}
\end{equation}
For each fixed $t<1$, $\mathcal M(\nu,1-t^2)\to1$ as $\nu\to\infty$.
Across the full scalar domain, the limiting margin vanishes only at $t=t_*$.
This point lies in the already proved range $t<1/\sqrt{10}$, not in the
closure of the remaining strip~\eqref{eq:scalar-atom-strip}; the latter has
positive limiting margin at least
$\ke(1/\sqrt{10}-t_*)^2$.
\end{proposition}
\end{samepage}

\begin{proof}
Put $c=1-t^2$ and
$U(\theta)=(1-c\sin^2(\theta/2))^{1/2}$, with normalized Lebesgue
measure on $[0,\pi]$. If $1\leq\nu\leq2$, Lyapunov's inequality gives
\[
 \E U^\nu\leq(\E U^2)^{\nu/2}
 =(1-c/2)^{\nu/2}.
\]
Maximizing the resulting squared bound over $c\in[0,1]$ gives
\[
 \mathcal M(\nu,c)^2
 \leq\frac{\pi\nu}{\nu+1}
       \left(\frac{\nu}{\nu+1}\right)^\nu
 \leq\frac{8\pi}{27}<1.
\]
The middle expression is increasing on $[1,2]$, and $\pi<27/8$ proves the
last inequality. This proves the first range.

Suppose next that $9/10\leq c<1$. For $x\in[0,\pi/2]$, set
\[
 D_c(x)=-\log(1-c\sin^2x)-cx^2.
\]
The derivative of
$H_c(x)=\sin(2x)/(1-c\sin^2x)-2x$ is
\[
 H_c'(x)=
 -\frac{2\sin^2x\{c^2\sin^2x-3c+2\}}
        {(1-c\sin^2x)^2}.
\]
The factor in braces changes sign once. Since $H_c(0)=0$ and
$H_c(\pi/2)=-\pi$, the function $D_c$ first increases and then decreases,
so its minimum occurs at an endpoint. The endpoint $x=0$ gives zero, while
\[
 D_c(\pi/2)=-\log(1-c)-\frac{c\pi^2}{4}\geq0.
\]
It suffices to check $c=9/10$, using $\log10>23/10$ and $\pi<22/7$;
the endpoint expression then increases with $c$. Therefore
$1-c\sin^2(\theta/2)<e^{-c\theta^2/4}$ when $0<\theta\leq\pi$.
At $c=1$, the same strict comparison follows from $\tan x>x$ for
$0<x<\pi/2$.
Integrating over the original finite interval gives
\[
 \mathcal M(\nu,c)
 <\operatorname{erf}\!\left(\pi\sqrt{\frac{\nu c}{8}}\right)<1
 \qquad(\nu>0).
\]
This proves the second range.

For the third range put $W=\nu c/4$. Equation~\eqref{eq:bessel-majorant} gives
\[
 \mathcal M(\nu,c)
 \leq e^{-W}I_0(W)\sqrt{2\pi W}<\frac65.
\]
For $0<W\leq1$, the series inequality $I_0(W)\leq e^{W^2/4}$ reduces the
claim to the increasing function
$\sqrt{2\pi W}\,e^{-W+W^2/4}$ and its value at one. For $W\geq1$, let
$r=I_1/I_0$. Termwise comparison gives $r(1)<1/2$. The Riccati identity
$r'=1-r^2-r/W$ then prevents $r$ from crossing
$1-1/(2W)$, because at a contact its derivative is $1/(4W^2)$ whereas
the comparison function has derivative $1/(2W^2)$. Thus the displayed
Bessel expression decreases for $W\geq1$. Finally,
$\sqrt{2\pi}e^{-3/4}<6/5$ follows from $\pi<22/7$ and the first five
positive terms of $e^{3/2}$. When $t\geq t_0$, the right side of the
scalar inequality~\eqref{eq:signed-atom-scalar} is at least $6/5$.

For fixed $c>0$, the final limit follows from the change of variables
$\theta=\sqrt{8/(\nu c)}\,u$ and dominated convergence in the resulting
Laplace integral. The limiting right side is $1+\ke(t-t_*)^2$, whose
gap from one vanishes only at $t=t_*$.
\end{proof}

\begin{samepage}
\begin{proposition}[Full scalar atom bound]\label{prop:scalar-atom-bound}
The scalar inequality~\eqref{eq:signed-atom-scalar} holds for every
$\nu\geq1$ and $0\leq t<1$, strictly at every finite point. Its margin can
tend to zero only when $\nu\to\infty$ and $t\to t_*$.
\end{proposition}
\end{samepage}

\begin{proof}
Proposition~\ref{prop:scalar-atom-ranges} leaves only the strip in
the scalar region~\eqref{eq:scalar-atom-strip}. Put
\[
 a=\frac1{\sqrt{10}},\qquad c=1-t^2,\qquad
 J_\nu(t)=\frac2\pi\int_0^{\pi/2}
 (\cos^2\theta+t^2\sin^2\theta)^{\nu/2}\,d\theta.
\]
For integers $m\geq0$, beta integration gives
\begin{equation}
 q_m(t):=J_{2m}(t)
 =4^{-m}\sum_{k=0}^m
 \binom{2k}{k}\binom{2m-2k}{m-k}t^{2k}.              \label{eq:scalar-even-moment}
\end{equation}
If $2m-2\leq\nu\leq2m$, $m\geq2$, H\"older log-convexity and
$\lambda=(\nu-2m+2)/2$ give
\[
 J_\nu(t)\leq q_{m-1}(t)^{1-\lambda}q_m(t)^\lambda.
\]
Set $\ell_m(t)=\log(q_{m-1}(t)/q_m(t))>0$. Maximizing the resulting factor
$\sqrt\nu\exp\{-\nu\ell_m(t)/2\}$ over all positive real $\nu$, a larger
domain than the order slab, gives
\begin{equation}
 \mathcal M(\nu,1-t^2)
 \leq U_m(t):=
 \sqrt{\frac{\pi(1-t^2)}{2e\ell_m(t)}}\,
 q_{m-1}(t)e^{(m-1)\ell_m(t)}.                       \label{eq:scalar-slab-majorant}
\end{equation}

It remains to bound finitely many one-variable functions. Divide $[0,1]$
into the $2048$ exact rational intervals
$[j/2048,(j+1)/2048]$ and map them affinely onto $[a,t_0]$.
Outward Arb evaluation of the positive polynomial~\eqref{eq:scalar-even-moment}
and the slab majorant~\eqref{eq:scalar-slab-majorant} proves
\begin{equation}
 1+\ke(t-t_*)^2-U_m(t)>\frac{19}{1000}.              \label{eq:scalar-compact-gap}
\end{equation}
The compact gap~\eqref{eq:scalar-compact-gap} holds on every such interval
and for every $2\leq m\leq64$. These
$2048\cdot63=129024$ boxes cover
$2\leq\nu\leq128$ and $a\leq t\leq t_0$, including all four edges.
The recorded smallest outward lower endpoint exceeds $0.0194747706542$.
The compact interval proof and independent Arb implementation are, respectively,
\begin{flushleft}
\path{verify/finite/scalar_atom/compact_certificate.py}\\
\path{verify/finite/scalar_atom/independent_certificate.py}
\end{flushleft}
Their digests and replay commands are recorded in \path{certificates/finite/}.

For $\nu\geq128$, the Bessel majorant from the proof of
Proposition~\ref{prop:scalar-atom-ranges} is decreasing because
$W=\nu c/4\geq1$. Throughout the remaining strip,
\[
 W\geq32(1-t_0^2).
\]
Since $t\geq a>t_*$, the right side of the scalar inequality is at least
$1+\ke(a-t_*)^2$. Since the Bessel majorant decreases in $W$, its left side
is at most its value at $W=32(1-t_0^2)$. Thus one outward evaluation gives
\[
 1+\ke(a-t_*)^2
 -\sqrt{2\pi W}\,e^{-W}I_0(W)>0.0180176
 \qquad\bigl(W=32(1-t_0^2)\bigr).
\]
This proves the tail for every real order, and hence the full scalar
inequality. All finite-region bounds are strict. All regions outside the
small-skew range have the fixed positive margins just proved. In the
small-skew range, the preceding error-function estimate gives
\[
 1+\ke(t-t_*)^2-\mathcal M(\nu,1-t^2)
 >\ke(t-t_*)^2+
 \operatorname{erfc}\!\left(
   \pi\sqrt{\frac{\nu(1-t^2)}8}\right).
\]
Hence a vanishing-margin sequence must have $t\to t_*$ and
$\nu\to\infty$; the Laplace limit in
Proposition~\ref{prop:scalar-atom-ranges} shows that this boundary is attained
as a limit.
\end{proof}


\bibliographystyle{plainnat}
\bibliography{refs}
\end{document}
